% ============================================================
%  RAPPORT DE FIN D'ÉTUDES – ESPRIT
%  AI-Driven Credit Risk Prediction and Explainable
%  Decision Platform for Banking
%  Société d'accueil : MPSoft
% ============================================================


\documentclass[a4paper,12pt,oneside]{report}

% ── Encodage et langue ───────────────────────────────────────
\usepackage[utf8]{inputenc}
\usepackage[T1]{fontenc}
\usepackage[french]{babel}

% ── Mise en page ─────────────────────────────────────────────
\usepackage[top=2.5cm, bottom=2.5cm, left=3cm, right=2.5cm]{geometry}
\usepackage{setspace}
\onehalfspacing

% ── Typographie ──────────────────────────────────────────────
\usepackage{lmodern}
\usepackage{microtype}

% ── Couleurs (UNIQUEMENT rouge ESPRIT + gris ; AUCUN BLEU) ───
\usepackage[dvipsnames,svgnames,table]{xcolor}
\definecolor{EspritRed}{RGB}{180, 0, 0}
\definecolor{DarkGray}{RGB}{50, 50, 50}
\definecolor{LightGray}{RGB}{235, 235, 235}

% ── Graphiques ───────────────────────────────────────────────
\usepackage{graphicx}
\usepackage{float}                 % placement [H] strict des flottants
\usepackage{tikz}
\usepackage{pgffor}
\usetikzlibrary{arrows.meta, positioning, shapes.geometric, calc, fit, backgrounds}

% Dossiers où LaTeX cherche les images (Overleaf : téléverser les
% logos à la racine du projet ou dans images/ — rien d'autre à faire)
\graphicspath{{./}{images/}{figures/}}

% ── Images absentes : emplacement réservé automatique ────────
% Tant qu'un fichier image n'a pas été téléversé, un cadre vide
% portant son nom est affiché à sa place : le document compile
% sans erreur. Dès que le fichier est ajouté au projet, l'image
% s'affiche normalement, sans aucune modification du code.
%
% Pour un espace totalement vierge (sans cadre ni nom de fichier),
% remplacer \bfpmephframetrue par \bfpmephframefalse ci-dessous.
\usepackage{keyval}
\newif\ifbfpmephframe
\bfpmephframetrue

\makeatletter
\newlength{\bfpme@w}
\newlength{\bfpme@h}
% Clés acceptées par l'emplacement réservé (les autres sont ignorées)
\define@key{bfpmeph}{width}{\setlength{\bfpme@w}{#1}}
\define@key{bfpmeph}{height}{\setlength{\bfpme@h}{#1}}
\define@key{bfpmeph}{totalheight}{\setlength{\bfpme@h}{#1}}
\define@key{bfpmeph}{scale}{}
\define@key{bfpmeph}{angle}{}
\define@key{bfpmeph}{page}{}
\define@key{bfpmeph}{trim}{}
\define@key{bfpmeph}{clip}[true]{}
\define@key{bfpmeph}{draft}[true]{}
\define@key{bfpmeph}{keepaspectratio}[true]{}
\define@key{bfpmeph}{origin}{}
\define@key{bfpmeph}{viewport}{}
\define@key{bfpmeph}{bb}{}
\define@key{bfpmeph}{interpolate}[true]{}
\define@key{bfpmeph}{type}{}
\define@key{bfpmeph}{ext}{}

\newcommand{\bfpme@placeholder}[2]{%
  \begingroup
    \setlength{\bfpme@w}{0pt}\setlength{\bfpme@h}{0pt}%
    \setkeys{bfpmeph}{#1}%
    % dimensions par défaut si une seule (ou aucune) a été donnée
    \ifdim\bfpme@w=0pt \ifdim\bfpme@h=0pt
        \setlength{\bfpme@w}{0.6\linewidth}\setlength{\bfpme@h}{0.28\linewidth}%
      \else
        \setlength{\bfpme@w}{1.8\bfpme@h}%
      \fi
    \else
      \ifdim\bfpme@h=0pt \setlength{\bfpme@h}{0.45\bfpme@w}\fi
    \fi
    \ifbfpmephframe
      \fboxsep=0pt \fboxrule=0.4pt
      \textcolor{gray!55}{%
        \fbox{\parbox[c][\bfpme@h][c]{\bfpme@w}{\centering
          \tiny\textcolor{gray!75}{\texttt{\detokenize{#2}}}}}}%
    \else
      \rule{0pt}{\bfpme@h}\hspace{\bfpme@w}%
    \fi
  \endgroup
}

% Recherche du fichier : racine, images/ et figures/, avec ou sans
% extension (\IfFileExists n'explore pas \graphicspath tout seul).
\newif\ifbfpme@found
\newcommand{\bfpme@try}[1]{%
  \ifbfpme@found\else\IfFileExists{#1}{\bfpme@foundtrue}{}\fi}
\newcommand{\bfpme@lookup}[1]{%
  \bfpme@foundfalse
  \bfpme@try{#1}\bfpme@try{#1.pdf}\bfpme@try{#1.png}%
  \bfpme@try{#1.jpg}\bfpme@try{#1.jpeg}\bfpme@try{#1.eps}%
  \bfpme@try{images/#1}\bfpme@try{images/#1.pdf}\bfpme@try{images/#1.png}%
  \bfpme@try{images/#1.jpg}\bfpme@try{images/#1.jpeg}\bfpme@try{images/#1.eps}%
  \bfpme@try{figures/#1}\bfpme@try{figures/#1.pdf}\bfpme@try{figures/#1.png}%
  \bfpme@try{figures/#1.jpg}\bfpme@try{figures/#1.jpeg}\bfpme@try{figures/#1.eps}%
}

% Interception de \includegraphics : image réelle si le fichier
% existe, emplacement réservé sinon.
\let\bfpme@origincludegraphics\includegraphics
\renewcommand{\includegraphics}[2][]{%
  \bfpme@lookup{#2}%
  \ifbfpme@found
    \bfpme@origincludegraphics[#1]{#2}%
  \else
    \bfpme@placeholder{#1}{#2}%
  \fi
}
\makeatother

% ── Style des chapitres ──────────────────────────────────────
\usepackage{titlesec}

% Chapitres NUMÉROTÉS : double règle + "CHAPITRE N" + titre MAJUSCULES
\titleformat{\chapter}[display]
  {\normalfont\centering}
  {\rule{\linewidth}{2pt}\\[6pt]%
   {\large\bfseries CHAPITRE \thechapter}\\%
   \rule{\linewidth}{1pt}}
  {12pt}
  {\Huge\bfseries\MakeUppercase}
  [\vspace{4pt}\rule{\linewidth}{0pt}]
\titlespacing*{\chapter}{0pt}{-30pt}{40pt}

% Chapitres NON NUMÉROTÉS (numberless) : règle + titre MAJUSCULES
\titleformat{name=\chapter,numberless}[display]
  {\normalfont\centering}
  {}{0pt}
  {\rule{\linewidth}{2pt}\\[8pt]\Huge\bfseries\MakeUppercase}
  [\vspace{4pt}\rule{\linewidth}{0pt}]

% Sections / sous-sections : noir gras (pas de couleur)
\titleformat{\section}
  {\normalfont\Large\bfseries}{\thesection}{1em}{}
\titleformat{\subsection}
  {\normalfont\large\bfseries}{\thesubsection}{1em}{}
\titleformat{\subsubsection}
  {\normalfont\normalsize\bfseries}{\thesubsubsection}{1em}{}
\titlespacing*{\section}{0pt}{16pt}{8pt}
\titlespacing*{\subsection}{0pt}{14pt}{6pt}
\titlespacing*{\subsubsection}{0pt}{10pt}{4pt}

% Commande pour chapitres non numérotés (intro, conclusion, résumé, abstract)
\newcommand{\chapterunnumbered}[1]{%
  \chapter*{#1}%
  \addcontentsline{toc}{chapter}{#1}%
  \markboth{\MakeUppercase{#1}}{\MakeUppercase{#1}}%
}

% ── Table des matières ───────────────────────────────────────
\usepackage[titles]{tocloft}
\setlength{\cftbeforechapskip}{8pt}
\renewcommand{\cftchapfont}{\bfseries}
\renewcommand{\cftchappagefont}{\bfseries}
\renewcommand{\cftdotsep}{1}

% ── En-têtes et pieds de page ────────────────────────────────
% Aucun en-tête : ni titre de chapitre, ni filet. Seul le numéro
% de page figure en pied, centré.
\usepackage{fancyhdr}
\pagestyle{fancy}
\fancyhf{}
\fancyfoot[C]{\thepage}
\renewcommand{\headrulewidth}{0pt}
\renewcommand{\footrulewidth}{0pt}
% Les pages d'ouverture de chapitre suivent la même règle
\fancypagestyle{plain}{%
  \fancyhf{}%
  \fancyfoot[C]{\thepage}%
  \renewcommand{\headrulewidth}{0pt}%
}

% ── Lettrines ────────────────────────────────────────────────
\usepackage{lettrine}
\renewcommand{\LettrineFontHook}{\bfseries}

% ── Hyperliens (TOUT EN NOIR : aucun bleu) ───────────────────
\usepackage[
  colorlinks=true,
  linkcolor=black, citecolor=black, urlcolor=black,
  anchorcolor=black, filecolor=black, menucolor=black, runcolor=black,
  pdftitle={AI-Driven Credit Risk Prediction and Explainable Decision Platform for Banking},
  pdfauthor={[Votre Nom]}
]{hyperref}

% ── Tableaux et listes ───────────────────────────────────────
\usepackage{booktabs}
\usepackage{tabularx}
\usepackage{longtable}
\usepackage{multirow}
\usepackage{array}
\usepackage{enumitem}
\setlist[itemize,1]{label=\textbullet, leftmargin=2em}
\setlist[itemize,2]{label=---,         leftmargin=3em}
\newcolumntype{Y}{>{\raggedright\arraybackslash}X}

% ── Code ─────────────────────────────────────────────────────
\usepackage{listings}
\definecolor{CodeBg}{RGB}{248, 248, 248}
\definecolor{CodeRule}{RGB}{180, 0, 0}
\definecolor{CodeNum}{RGB}{150, 150, 150}
\definecolor{CodeComment}{RGB}{110, 110, 110}

% Style unifié : filet rouge ESPRIT à gauche, fond très clair, sans cadre lourd
\lstdefinestyle{bfpme}{
  basicstyle      = \ttfamily\footnotesize,
  backgroundcolor = \color{CodeBg},
  frame           = leftline,
  framerule       = 2pt,
  rulecolor       = \color{CodeRule},
  framexleftmargin  = 8pt,
  framexrightmargin = 4pt,
  framextopmargin   = 5pt,
  framexbottommargin= 5pt,
  xleftmargin     = 12pt,
  xrightmargin    = 2pt,
  numbers         = left,
  numberstyle     = \tiny\color{CodeNum},
  numbersep       = 9pt,
  keywordstyle    = \bfseries\color{DarkGray},
  commentstyle    = \itshape\color{CodeComment},
  stringstyle     = \color{DarkGray},
  breaklines      = true,
  breakatwhitespace = true,
  breakindent     = 1em,
  columns         = fullflexible,
  keepspaces      = true,
  showstringspaces= false,
  tabsize         = 2,
  aboveskip       = 12pt,
  belowskip       = 10pt,
  captionpos      = b,
  abovecaptionskip= 6pt
}
\lstset{style=bfpme}

% Environnement : un listing insécable, centré, jamais coupé entre deux pages
\newenvironment{codeblock}
  {\par\medskip\noindent\begin{minipage}{\linewidth}}
  {\end{minipage}\par\medskip}

% ── Mathématiques ────────────────────────────────────────────
\usepackage{amsmath, amssymb}

% ── Figures et légendes ──────────────────────────────────────
\usepackage{caption}
\usepackage{subcaption}
\captionsetup{labelsep=endash, font=small, labelfont=bf, justification=centering}

% ── Styles TikZ réutilisables ────────────────────────────────
\tikzset{
  layerbox/.style   = {draw=DarkGray, thick, rounded corners=3pt, fill=LightGray,
                       text width=11cm, minimum height=1.05cm, align=center, font=\small},
  flowbox/.style    = {draw=DarkGray, thick, rounded corners=2pt, fill=LightGray,
                       minimum height=0.9cm, minimum width=2.4cm, align=center, font=\footnotesize},
  stagebox/.style   = {draw=DarkGray, very thick, rounded corners=3pt, fill=LightGray,
                       text width=2.9cm, minimum height=1.6cm, align=center, font=\footnotesize},
  arrow/.style      = {-{Stealth[length=2.5mm]}, thick, DarkGray},
  % Styles du schéma de flux (fig. « Flux de données »).
  % NB : préfixe « flux » obligatoire — des noms courts comme « step »
  % entreraient en collision avec les clés réservées de TikZ.
  fluxbox/.style    = {draw=DarkGray, thick, rounded corners=2pt, fill=LightGray,
                       text width=2.3cm, minimum height=1.15cm, align=center,
                       font=\scriptsize},
  fluxio/.style     = {draw=DarkGray, very thick, rounded corners=6pt, fill=white,
                       text width=2.0cm, minimum height=1.15cm, align=center,
                       font=\scriptsize\bfseries},
  fluxlab/.style    = {font=\tiny\itshape, text=DarkGray, inner sep=1.5pt},
  fluxarrow/.style  = {-{Stealth[length=2.2mm]}, thick, DarkGray}
}

% ── Bibliographie ────────────────────────────────────────────
% Bibliographie native (thebibliography) : les citations \cite{...}
% s'affichent en [1], [2]... après deux passes de pdfLaTeX, SANS biber.
% Titre de la section :
\renewcommand{\bibname}{Webographie \& Bibliographie}

% ════════════════════════════════════════════════════════════
\begin{document}

% ════════════════════════════════════════════════════════════
%  PAGES PRÉLIMINAIRES  (numérotation arabe continue : 1, 2, 3...)
% ════════════════════════════════════════════════════════════
\pagenumbering{arabic}
\setcounter{page}{1}

\chapter*{Dédicace}
\thispagestyle{plain}

\vspace*{\fill}
\begin{center}
\itshape
À \textbf{Dieu Tout-Puissant}, pour Sa grâce et Ses bienfaits,\\[1.2em]

Je dédie ce modeste travail\,:\\[1.2em]

À mon très cher père, \textbf{Mongi}, ma source de force et mon modèle,\\[0.6em]
À ma très chère mère, \textbf{Dora}, pour son amour et ses sacrifices sans fin,\\[0.6em]
À mes chères sœurs, \textbf{Sarah} et \textbf{Heger}, pour leur soutien et leur tendresse,\\[0.6em]
À ma grand-mère, \textbf{Beya}, pour ses prières et sa bienveillance,\\[0.6em]
À mes amis, \textbf{Sarra}, \textbf{Syrine}, \textbf{Dhia} et \textbf{Ramy}, pour leur présence et leurs encouragements.\\[1.6em]

En témoignage de ma profonde reconnaissance et de mon affection.
\end{center}
\vspace*{\fill}

\chapter*{Remerciements}
\thispagestyle{plain}

Au terme de ce projet de fin d'études, je tiens à exprimer ma profonde gratitude à toutes
les personnes qui ont contribué, de près ou de loin, à son aboutissement.

Je remercie tout particulièrement mes encadrants professionnels, \textbf{Monsieur Lassaad Maalej}
et \textbf{Madame Ikbel Aloui}, ainsi que mon encadrante académique, \textbf{Madame Saoussen
Lakhdher}, pour leur disponibilité, leurs conseils précieux et la confiance qu'ils m'ont accordée
tout au long de ce travail. Leur accompagnement et leur expertise ont été déterminants pour la
réussite de ce projet.

Mes remerciements s'adressent également à l'ensemble de l'équipe de \textbf{MPSoft} pour son
accueil chaleureux et son esprit de collaboration, ainsi qu'au corps enseignant de l'\textbf{ESPRIT}
pour la qualité de la formation dispensée.

Enfin, j'exprime toute ma reconnaissance aux membres du jury pour l'honneur qu'ils me font en
acceptant d'évaluer ce travail.

{
  \renewcommand{\contentsname}{Table des matières}
  \tableofcontents
}

\chapter*{Liste des abréviations}
\thispagestyle{plain}
\begin{tabular}{ll}
  \textbf{IA}    & Intelligence Artificielle \\
  \textbf{ML}    & Machine Learning (apprentissage automatique) \\
  \textbf{LLM}   & Large Language Model (grand modèle de langage) \\
  \textbf{XAI}   & Explainable Artificial Intelligence \\
  \textbf{SHAP}  & SHapley Additive exPlanations \\
  \textbf{PD}    & Probabilité de Défaut (\textit{Probability of Default}) \\
  \textbf{AUC}   & Aire sous la courbe ROC (\textit{Area Under the Curve}) \\
  \textbf{ROC}   & Receiver Operating Characteristic \\
  \textbf{VaR}   & Value-at-Risk \\
  \textbf{EDA}   & Analyse exploratoire des données (\textit{Exploratory Data Analysis}) \\
  \textbf{PSI}   & Population Stability Index \\
  \textbf{NLP}   & Traitement du Langage Naturel \\
  \textbf{API}   & Interface de Programmation Applicative \\
  \textbf{REST}  & Representational State Transfer \\
  \textbf{SCRUM} & Cadre de gestion de projet Agile \\
  \textbf{UML}   & Langage de Modélisation Unifié \\
  \textbf{PME}   & Petite et Moyenne Entreprise \\
\end{tabular}

% ════════════════════════════════════════════════════════════
%  INTRODUCTION GÉNÉRALE
% ════════════════════════════════════════════════════════════
\clearpage

\chapterunnumbered{Introduction générale}

\lettrine[lines=3, lhang=0.1, loversize=0.05]{D}{}ans un secteur bancaire en pleine
mutation, où la digitalisation transforme profondément les pratiques du crédit, l'évaluation
du risque de défaut demeure un enjeu stratégique majeur. Les établissements financiers
font face à une double exigence : prendre des décisions rapides et fiables, et garantir la
transparence et l'explicabilité de chaque décision vis-à-vis des régulateurs, des auditeurs
et des clients.

Les approches traditionnelles de scoring de crédit, fondées sur des modèles statistiques
linéaires et des règles expertes, atteignent aujourd'hui leurs limites. Elles peinent à
capturer la complexité des comportements de remboursement, à s'adapter à des profils
d'emprunteurs hétérogènes, et à fournir des justifications compréhensibles pour chaque
décision individuelle. Par ailleurs, la «~boîte noire~» inhérente aux modèles d'apprentissage
automatique avancés constitue un frein réglementaire et éthique majeur à leur adoption
dans un contexte bancaire.

L'intelligence artificielle explicable (\textit{Explainable AI}, XAI) ouvre une voie nouvelle
en conciliant la puissance prédictive des algorithmes modernes et l'exigence d'auditabilité.
En parallèle, l'émergence des agents IA et des grands modèles de langage (LLM) offre des
perspectives inédites pour orchestrer des processus décisionnels complexes, interagir en
langage naturel avec les analystes, et produire des explications contextualisées et accessibles.

C'est dans ce cadre que s'inscrit notre projet de fin d'études, réalisé au sein de
Manager Partner Software (MPSoft), leader tunisien en solutions bancaires. Notre mission
consiste à concevoir et développer une plateforme d'aide à la décision pour le risque de
crédit, pilotée par un agent IA, combinant la modélisation par apprentissage automatique de
la probabilité de défaut (PD), l'explicabilité globale et locale par SHAP, le raisonnement en
langage naturel via un LLM et la simulation de scénarios de stress par méthode de Monte Carlo.
L'ensemble de la plateforme intègre une logique de supervision humaine
(\textit{human-in-the-loop}) garantissant que l'IA assiste~--- sans jamais remplacer~--- le
jugement des analystes crédit et des comités de décision des établissements bancaires.

Ce rapport est organisé en sept chapitres suivant une progression logique, de la
compréhension métier jusqu'à la validation du système livré.

Le Chapitre~1 établit le cadre général du projet : il présente la société
d'accueil MPSoft, analyse le processus existant d'évaluation du risque
de crédit PME, formule la problématique, décrit la solution proposée, justifie le choix
méthodologique SCRUM et spécifie les besoins fonctionnels et non fonctionnels.

Le Chapitre~2 dresse un état de l'art couvrant le risque de crédit bancaire et
ses spécificités PME, les approches traditionnelles et modernes de scoring, les techniques
d'apprentissage automatique, les techniques d'explicabilité (XAI) telles que SHAP, les
architectures d'agents IA, l'usage des grands modèles de langage en finance et les méthodes
de simulation Monte Carlo.

Le Chapitre~3 détaille la conception et l'architecture technique de la plateforme :
architecture multi-couches, conception des données, architecture de l'agent IA, diagrammes
UML, environnement de développement et planification des sprints et des releases.

Le Chapitre~4 (Release~I) couvre la compréhension métier du crédit PME et
l'ingénierie des données : cycle de vie du crédit, collecte et exploration des données
bancaires, nettoyage, \textit{feature engineering} et évaluation de la qualité et des biais.

Le Chapitre~5 (Release~II) présente la modélisation prédictive (régression
logistique, arbre de décision, Random Forest, Gradient Boosting) ainsi que l'explicabilité
(SHAP global et local, auditabilité des décisions).

Le Chapitre~6 (Release~III) décrit l'agent IA et son orchestration décisionnelle,
ainsi que l'intégration d'un LLM pour la génération d'explications en langage naturel et
l'interaction avec les analystes.

Le Chapitre~7 (Release~IV) aborde les tests de stress du portefeuille par
simulation Monte Carlo : construction de scénarios adverses, estimation de la distribution
des pertes et des indicateurs de risque extrême du portefeuille PME.

La conclusion générale synthétise les apports du projet et ouvre sur les
perspectives d'évolution de la plateforme.

% ════════════════════════════════════════════════════════════
%  CHAPITRE 1 – CONTEXTE GÉNÉRAL ET SPÉCIFICATION DES BESOINS
% ════════════════════════════════════════════════════════════
\chapter{Contexte général et spécification des besoins}

\section*{Introduction}

Dans un contexte d'accélération numérique où la transformation digitale est passée du
statut de concept à celui de nécessité opérationnelle, les institutions financières font
face à un défi de taille~: évaluer le risque de crédit de manière précise, rapide et
explicable. Ce constat est particulièrement critique pour les banques spécialisées dans le
financement des PME, dont les portefeuilles présentent des caractéristiques de risque
spécifiques.

Ce chapitre présente le cadre général de notre projet de fin d'études, réalisé au sein de
MPSoft, leader tunisien en solutions bancaires. Notre travail
s'articule autour du développement d'une plateforme d'aide à la décision pour le risque de
crédit, combinant la modélisation par apprentissage automatique de la probabilité de défaut
(PD), l'intelligence artificielle explicable, un agent orchestrateur et le raisonnement LLM.

\section{Cadre général du projet}

Le travail présenté dans ce rapport a été effectué au sein de la société Manager Partner
Software (MPSoft) dans le cadre d'un projet de fin d'études en vue de l'obtention du
diplôme national d'ingénieur en Informatique. Ce projet vise à moderniser les outils
d'évaluation du risque de crédit des établissements bancaires et à résoudre des
problématiques concrètes d'aide à la décision dans le contexte du financement des PME.

\subsection{Objectifs du projet}

L'objectif principal de ce projet est de concevoir et développer un agent IA
d'aide à la décision pour le risque de crédit des PME, alliant haute performance prédictive
et pleine explicabilité des décisions. Cet agent doit~:

\begin{itemize}
  \item prédire la probabilité de défaut (PD) de toute PME emprunteuse à partir de ses
        caractéristiques financières, sectorielles et comportementales~;
  \item expliquer chaque décision de crédit à travers des valeurs SHAP globales et locales,
        compréhensibles tant par les analystes de la banque que par les régulateurs~;
  \item permettre aux analystes crédit d'interroger le système en langage naturel sur les
        facteurs de risque et les résultats~;
  \item tester la robustesse du portefeuille PME sous des scénarios de stress économique
        par simulation Monte Carlo.
\end{itemize}

\section{Organisme d'accueil}

\subsection{Présentation de MPSoft}

\begin{figure}[H]
  \centering
  \definecolor{MPGray}{RGB}{176,176,176}
  \definecolor{MPRed}{RGB}{213,20,32}
  \resizebox{4.6cm}{!}{%
  \begin{tikzpicture}
    \fill[MPGray, rounded corners=2pt] (0,0) rectangle (4,3.4);
    \node[font=\fontsize{48}{48}\bfseries\sffamily, text=white] at (2,2.45) {MP};
    \node[font=\fontsize{42}{42}\bfseries\sffamily, text=MPRed]  at (2,1.05) {Soft};
  \end{tikzpicture}}\\[0.18cm]
  {\footnotesize\sffamily\color{gray!70}%
    M\,A\,N\,A\,G\,E\,R~~P\,A\,R\,T\,N\,E\,R~~S\,O\,F\,T\,W\,A\,R\,E}
  \caption{Logo de Manager Partner Software}
  \label{fig:logo_mpsoft}
\end{figure}

MPSoft constitue l'environnement professionnel dans lequel s'inscrit ce projet de fin
d'études. Fondée en 2001, cette entreprise tunisienne s'est positionnée comme un
leader tunisien en solutions bancaires, développant une expertise reconnue avec
plus de 40~modules métiers (MpBank, MpPost) et une plateforme No-Code révolutionnaire.

\subsection{Mission et Valeurs}

MPSoft révolutionne le développement des systèmes d'information bancaires en offrant des
solutions modulaires, performantes, sécurisées et conformes aux exigences réglementaires.
Pour concrétiser cette ambition, l'entreprise s'investit dans~:

\begin{itemize}
  \item systèmes bancaires modulaires et personnalisés : solutions adaptées
        pour banques de détail, banques postales, microfinance ou finance islamique~;
  \item déploiement rapide avec ROI optimisé : technologies No-Code et
        architectures modulaires réduisant les délais jusqu'à 5x~;
  \item performance, conformité et adaptabilité locale : solutions compatibles
        Open API et ajustées aux réglementations locales~;
  \item partenariat sur le long terme : accompagnement des institutions tout au
        long de leur digitalisation~;
  \item accessibilité accrue : plateforme No-Code permettant aux profils non
        techniques de concevoir des applications métiers.
\end{itemize}

\subsection{Organigramme de Manager Partner Software}

\begin{figure}[H]
  \centering
  \definecolor{OrgCrimson}{RGB}{168,24,58}
  \definecolor{OrgOrange}{RGB}{240,130,20}
  \definecolor{OrgBlue}{RGB}{20,70,150}
  \definecolor{OrgLBlue}{RGB}{150,200,235}
  \resizebox{\textwidth}{!}{%
  \begin{tikzpicture}[
      box/.style={rounded corners=1pt, minimum width=3.4cm, minimum height=0.95cm,
                  align=center, text=white, font=\bfseries, inner sep=3pt},
      lo/.style={draw=OrgOrange, line width=1.3pt, rounded corners=6pt},
      lb/.style={draw=OrgBlue,   line width=1.3pt, rounded corners=6pt},
      ll/.style={draw=OrgLBlue,  line width=1.3pt, rounded corners=6pt}]

    % ── Direction générale ──
    \node[fill=OrgCrimson, text=white, minimum width=1.15cm, minimum height=5.6cm,
          font=\Large\bfseries, rounded corners=1pt] (dg) at (0,0) {\rotatebox{90}{DG}};

    % ── Niveau 1 : directions (orange) ──
    \node[box, fill=OrgOrange] (mc)   at (4.1, 1.95) {Marketing+Commercial};
    \node[box, fill=OrgOrange] (af)   at (4.1, 0.00) {Administratif+Financier};
    \node[box, fill=OrgOrange] (tech) at (4.1,-1.95) {Technique};

    % ── Niveau 2 : pôles techniques (bleu) ──
    \node[box, fill=OrgBlue] (infra)  at (8.6, 0.40) {Infrastructure};
    \node[box, fill=OrgBlue] (dev)    at (8.6,-1.35) {Développement};
    \node[box, fill=OrgBlue] (met)    at (8.6,-3.10) {Métier};

    % ── Niveau 3 : sous-pôles infrastructure (bleu clair) ──
    \node[box, fill=OrgLBlue] (sec)   at (13.1, 1.75) {Sécurité};
    \node[box, fill=OrgLBlue] (res)   at (13.1, 0.40) {Réseau};
    \node[box, fill=OrgLBlue] (cloud) at (13.1,-0.95) {Cloud et DeVops};

    % ── Connexions DG -> directions ──
    \coordinate (t1) at ($(dg.east)+(0.55,0)$);
    \draw[lo] (dg.east)   -- (t1);
    \draw[lo] (t1|-mc)    -- (t1|-tech);
    \draw[lo] (t1|-mc)    -- (mc.west);
    \draw[lo] (t1|-af)    -- (af.west);
    \draw[lo] (t1|-tech)  -- (tech.west);

    % ── Connexions Technique -> pôles ──
    \coordinate (t2) at ($(tech.east)+(0.55,0)$);
    \draw[lb] (tech.east) -- (t2|-tech);
    \draw[lb] (t2|-infra) -- (t2|-met);
    \draw[lb] (t2|-infra) -- (infra.west);
    \draw[lb] (t2|-dev)   -- (dev.west);
    \draw[lb] (t2|-met)   -- (met.west);

    % ── Connexions Infrastructure -> sous-pôles ──
    \coordinate (t3) at ($(infra.east)+(0.55,0)$);
    \draw[ll] (infra.east) -- (t3|-infra);
    \draw[ll] (t3|-sec)    -- (t3|-cloud);
    \draw[ll] (t3|-sec)    -- (sec.west);
    \draw[ll] (t3|-res)    -- (res.west);
    \draw[ll] (t3|-cloud)  -- (cloud.west);
  \end{tikzpicture}}
  \caption{Organigramme de Manager Partner Software}
  \label{fig:orga_mpsoft}
\end{figure}

\subsection{Cadre du stage}

Le stage s'est déroulé au sein de l'équipe Data \& Innovation de MPSoft. L'objectif consiste
à doter les établissements bancaires d'un outil d'évaluation et d'aide à la décision crédit
de nouvelle génération, capable de répondre aux spécificités du risque PME.

\section{Étude de l'existant et problématique}

\subsection{Le financement des PME par les banques}

Les banques de financement des petites et moyennes entreprises jouent un rôle clé dans
le soutien au développement du tissu entrepreneurial. Leur mission consiste à financer la
création, l'extension et la modernisation des PME, en accordant des crédits à moyen et long
terme à des projets dont les profils de risque sont souvent éloignés des standards de la banque de
détail classique. Les portefeuilles de crédit PME se caractérisent par~:

\begin{itemize}
  \item des volumes unitaires de crédit élevés comparés au retail, mais des
        échantillons globaux relativement limités~;
  \item une forte hétérogénéité sectorielle (industrie, services, agriculture,
        TIC, tourisme)~;
  \item une prédominance du risque idiosyncrasique lié au projet, au promoteur
        et à la qualité de gestion de la PME~;
  \item une forte sensibilité aux cycles macro-économiques et aux chocs
        sectoriels~;
  \item un horizon de remboursement long (5 à 12 ans), accroissant l'incertitude.
\end{itemize}

\subsection{Analyse de l'existant}

L'analyse approfondie des processus actuels d'évaluation du risque de crédit au sein des banques
révèle une situation contrastée~:

\begin{itemize}
  \item Points forts existants :
  \begin{itemize}
    \item expertise métier solide des analystes crédit sur les segments PME~;
    \item comité de crédit structuré avec une gouvernance claire~;
    \item connaissance fine des spécificités sectorielles du tissu PME~;
    \item historique de données client conséquent.
  \end{itemize}
  \item Limitations identifiées :
  \begin{itemize}
    \item scoring largement manuel, reposant sur des grilles expertes statiques~;
    \item faible pouvoir prédictif sur les profils complexes ou les secteurs émergents~;
    \item absence d'explicabilité formalisée pour justifier un refus individuel~;
    \item aucun mécanisme automatisé de test de stress~;
    \item délais d'instruction longs liés à la nature manuelle du processus.
  \end{itemize}
\end{itemize}

\subsection{Problématique}

\subsubsection{Limites des systèmes d'évaluation actuels}

Les outils de scoring actuellement utilisés par les banques reposent sur des grilles expertes
linéaires dont la capacité de discrimination plafonne sur les profils PME les plus
complexes. La diversité des secteurs financés et l'hétérogénéité des structures financières
des PME exigent des modèles capables de capturer des interactions non linéaires
entre variables financières, sectorielles et comportementales.

\subsubsection{Enjeux réglementaires et explicabilité}

La réglementation impose une traçabilité complète des décisions de crédit. Tout refus doit
pouvoir être justifié de manière individuelle, compréhensible et auditable. Cette exigence
d'explicabilité, aujourd'hui absente de nombreux outils en production dans les banques, est au
c\oe ur de notre démarche.

\subsubsection{Conséquences opérationnelles}

Ces lacunes se traduisent concrètement par des délais d'instruction allongés, un taux
d'erreur non négligeable sur les décisions de crédit, et une exposition mal maîtrisée du
portefeuille PME aux risques de stress macro-économique et sectoriel.

\section{Solution proposée}

\subsection{Modélisation par apprentissage automatique}

Le c\oe ur prédictif de la plateforme repose sur des modèles d'apprentissage automatique
(\textit{machine learning}) entraînés à estimer la probabilité de défaut d'une PME à partir de
ses caractéristiques financières, sectorielles et comportementales. La démarche de modélisation
suit un protocole rigoureux~: constitution et nettoyage d'un jeu de données, ingénierie de
variables discriminantes (ratios de solvabilité, indices de pression de paiement), puis
entraînement et comparaison de plusieurs familles d'algorithmes. Des modèles de référence
interprétables (régression logistique, arbre de décision) fixent une base de comparaison,
tandis que des modèles ensemblistes plus puissants (Random Forest, Gradient Boosting via
XGBoost) capturent les interactions non linéaires propres au risque PME. La sélection du modèle
final résulte d'un arbitrage entre pouvoir discriminant (AUC-ROC), stabilité et
interprétabilité, évalué par validation croisée. Le préprocesseur et le modèle retenu sont
sérialisés dans un artefact unique, garantissant la reproductibilité entre l'entraînement et
l'inférence en production.

\subsection{Module d'explicabilité (XAI)}

Le module XAI produit, pour chaque demande de crédit PME, des explications globales
(importance des variables à l'échelle du portefeuille) et locales (contribution SHAP par
variable pour la PME analysée), dans un format auditable par les analystes de la banque et les
régulateurs.

\subsection{Agent IA pour la décision de crédit}

Nous proposons un agent IA orchestrateur qui pilote l'ensemble du pipeline
décisionnel bancaire~: collecte et traitement des données PME, exécution du modèle
prédictif, application des seuils de confiance, escalade humaine vers le comité de crédit
si nécessaire, et journalisation de chaque décision.

\subsection{Interaction en langage naturel (LLM)}

Un grand modèle de langage (LLM) est intégré comme couche de raisonnement et de
communication~: il génère des explications narratives à partir des sorties du modèle, et
permet aux analystes crédit d'interroger le système en langage naturel sur les
facteurs de risque et les scénarios.

\subsection{Tests de stress par simulation Monte Carlo}

La plateforme intègre un moteur de simulation Monte Carlo pour tester la robustesse du
portefeuille PME sous des scénarios adverses : à partir d'un dossier de base et d'un ensemble
de chocs appliqués aux variables clés, le moteur génère un grand nombre de scénarios perturbés
et reconstitue la distribution des probabilités de défaut, dont sont dérivés des indicateurs de
risque extrême (VaR, Expected Shortfall).

\subsection{Valeur ajoutée}

Cette solution transforme un processus expert de plusieurs jours en une analyse automatisée,
explicable et robuste, réalisée en quelques minutes~; elle autonomise les équipes
d'analystes sans pour autant automatiser les décisions finales du comité de crédit.

\section{Choix méthodologique}

\subsection{Étude comparative des méthodes Agile}

Plusieurs cadres agiles ont été comparés (Scrum, Kanban, XP) selon des critères de
structuration, d'itérativité et d'adaptation aux retours métier, afin d'identifier la
méthode la plus adaptée à un projet data science incrémental.

\subsection{Justification du choix de SCRUM}

La méthode SCRUM a été retenue pour sa capacité à structurer le développement en sprints
itératifs, permettant une validation progressive des livrables par les parties prenantes
de MPSoft et les analystes crédit des banques, et une adaptation continue aux retours terrain.

\subsection{Équipe SCRUM}

L'équipe projet se répartit selon les rôles SCRUM classiques : un Product Owner garant de
la valeur métier, un Scrum Master facilitateur, et une équipe de développement chargée de
la réalisation technique.

\subsection{Formalisme de modélisation}

Nous utilisons le langage UML pour modéliser les cas d'utilisation, les séquences
d'interaction, le modèle de classes et l'architecture de déploiement de la plateforme.

\subsection{Événements SCRUM implémentés}

Les cérémonies SCRUM mises en œuvre comprennent la planification de sprint, la mêlée
quotidienne, la revue de sprint et la rétrospective, garantissant un rythme de livraison
régulier et une amélioration continue.

\section{Spécification des besoins}

\subsection{Besoins fonctionnels}

\subsubsection{Prédiction du risque de crédit PME}

Le système doit estimer la probabilité de défaut de chaque PME à partir de ses
caractéristiques financières, sectorielles et comportementales.

\subsubsection{Explicabilité et traçabilité des décisions (XAI)}

Chaque prédiction doit être accompagnée d'explications globales et locales auditables, et
journalisée pour assurer la traçabilité.

\subsubsection{Interaction en langage naturel avec les analystes}

Les analystes doivent pouvoir interroger le système en langage naturel et obtenir des
réponses contextualisées sur les facteurs de risque.

\subsubsection{Tests de stress et simulation Monte Carlo}

Le système doit permettre de simuler des scénarios adverses et d'évaluer la robustesse du
portefeuille.

\subsection{Besoins non fonctionnels}

\subsubsection{Performance et scalabilité}

La plateforme doit traiter les dossiers dans des délais courts et supporter la montée en
charge du portefeuille.

\subsubsection{Sécurité, confidentialité et protection des données}

Les données sensibles des PME doivent être protégées conformément aux bonnes pratiques de
sécurité.

\subsubsection{Auditabilité et gouvernance}

L'ensemble des décisions et traitements doit être traçable et gouverné, de manière à
satisfaire les exigences d'audit interne et externe.

\subsubsection{Maintenabilité et évolutivité}

L'architecture modulaire doit faciliter la maintenance et l'ajout de nouvelles
fonctionnalités.

\section*{Conclusion}

Ce premier chapitre a posé les fondations du travail~: la présentation de MPSoft (société
d'accueil), l'analyse critique du processus existant, la définition
de la problématique et la description de la solution proposée. Le chapitre suivant dresse
un état de l'art approfondi formant le socle scientifique et technique de la plateforme.

% ════════════════════════════════════════════════════════════
%  CHAPITRE 2 – ÉTAT DE L'ART
% ════════════════════════════════════════════════════════════
\chapter{État de l'art}

\section*{Introduction}

Ce chapitre dresse un panorama des concepts, techniques et cadres méthodologiques mobilisés
dans notre projet : le risque de crédit et ses spécificités PME, les approches traditionnelles
de scoring, l'apprentissage automatique appliqué au crédit, l'explicabilité (XAI), les agents
IA, les LLM en finance et la simulation Monte Carlo. Pour chaque thème, nous positionnons les
choix retenus par rapport à l'état de la recherche et aux pratiques de l'industrie.

\section{Le risque de crédit bancaire}

Le risque de crédit constitue la principale source de risque pour une banque de financement.
Cette section en présente les concepts fondamentaux.

\subsection{Définition et cycle de vie du crédit}

Le risque de crédit désigne la perte potentielle résultant de l'incapacité d'un emprunteur
à honorer ses engagements~\cite{thomas2002}. Le cycle de vie d'un crédit se décompose en quatre
phases : la demande et l'instruction, l'analyse et l'approbation, le décaissement et le suivi,
puis, le cas échéant, le défaut et le recouvrement. Chacune génère des données exploitables
pour l'évaluation du risque, depuis les états financiers collectés à l'instruction jusqu'aux
signaux comportementaux observés durant le suivi.

\subsection{Catégories de risque : éligibilité et probabilité de défaut}

On distingue le risque d'éligibilité, qui détermine si un dossier remplit les conditions
minimales d'octroi, de la probabilité de défaut (PD), qui quantifie la probabilité qu'un
emprunteur fasse défaut sur un horizon donné. La PD, généralement estimée sur un horizon de
douze mois, constitue le paramètre central de notre plateforme et la cible de la modélisation
prédictive.

\subsection{Spécificités du risque PME}

Le risque associé aux PME se distingue par sa nature fortement idiosyncrasique : il dépend
du projet, de la qualité du promoteur et de la structure financière de l'entreprise. À la
différence du crédit aux particuliers, les volumes d'observations sont plus limités et la
sensibilité aux chocs sectoriels et macro-économiques est plus marquée. Ces caractéristiques
compliquent l'estimation statistique du risque : la rareté relative des événements de défaut
et l'hétérogénéité des profils appellent des modèles capables de tirer parti d'un signal
dispersé sur de nombreuses variables, tout en restant robustes au sur-apprentissage.

\subsection{Rôle des parties prenantes : analystes, gestionnaires, comités}

L'évaluation du crédit mobilise plusieurs acteurs : les analystes instruisent les dossiers,
les gestionnaires de risque surveillent l'exposition globale, et les comités de crédit
prennent la décision finale. Notre plateforme vise à assister ces acteurs sans se
substituer à leur jugement, dans une logique de supervision humaine.

\section{Approches traditionnelles de scoring de crédit}

Les approches traditionnelles reposent sur des modèles statistiques linéaires et des grilles
expertes, dont l'histoire remonte aux travaux fondateurs sur la prédiction de la faillite
d'entreprise~\cite{altman1968}.

\subsection{Modèles statistiques classiques}

Les modèles classiques s'appuient principalement sur la régression logistique et l'analyse
discriminante, qui estiment une probabilité de défaut à partir d'une combinaison linéaire de
variables~\cite{hosmer2013}. La régression logistique, en particulier, demeure la référence
réglementaire du \textit{credit scoring} : elle fournit des coefficients directement
interprétables sous forme de rapports de cotes (\textit{odds ratios}), ce qui facilite la
justification des décisions auprès des régulateurs. Son interprétabilité élevée en a fait le
standard pendant plusieurs décennies.

\subsection{Scorecards expertes}

Les scorecards expertes traduisent l'expérience métier en grilles de notation pondérées,
attribuant des points selon les caractéristiques du dossier. Largement utilisées en
financement PME, elles offrent transparence et simplicité mais peinent à capturer les
interactions complexes entre variables. Le moteur de scoring interne BFPME, dont s'inspire la
génération de notre jeu de données, relève précisément de cette catégorie.

\subsection{Limites des approches existantes}

Ces approches souffrent de leur hypothèse de linéarité, de leur faible capacité à modéliser
les interactions, et d'un manque d'adaptation aux profils atypiques. Elles ne fournissent
aucune explication locale formalisée par décision, ce qui limite leur auditabilité au niveau
individuel. Les études comparatives de grande ampleur montrent que, sur données de crédit, les
méthodes d'apprentissage automatique modernes surpassent significativement ces approches
linéaires en pouvoir de discrimination~\cite{lessmann2015}.

\section{Apprentissage automatique appliqué au risque de crédit}

Les modèles modernes de machine learning offrent une capacité prédictive supérieure en
capturant les interactions non linéaires entre variables, au prix d'une interprétabilité
directe plus faible qu'il convient de compenser par des méthodes d'explicabilité dédiées.

\subsection{Arbres de décision et modèles ensemblistes}

L'arbre de décision partitionne récursivement l'espace des variables en régions homogènes
au regard de la cible ; simple et lisible, il souffre néanmoins d'une forte variance. Les
méthodes ensemblistes remédient à cette limite en combinant de multiples arbres. Le Random
Forest~\cite{breiman2001} agrège des arbres entraînés sur des sous-échantillons aléatoires
(\textit{bagging}) et réduit ainsi la variance, tandis que le Gradient
Boosting~\cite{friedman2001} construit les arbres séquentiellement, chaque arbre corrigeant les
erreurs résiduelles du précédent. Ces modèles capturent efficacement les relations non
linéaires propres au risque PME.

\subsection{Le Gradient Boosting et XGBoost}

Parmi les implémentations de Gradient Boosting, XGBoost~\cite{chen2016} s'est imposé comme une
référence en compétition et en production grâce à une optimisation régularisée, une gestion
native des valeurs manquantes et une parallélisation efficace. Sa robustesse et sa performance
sur des données tabulaires hétérogènes en font un candidat naturel pour le scoring de crédit,
et il constitue le modèle déployé au c\oe ur de notre plateforme~\cite{xgboostdoc}.

\subsection{Étude comparative des approches ML}

Le choix d'un modèle résulte d'un arbitrage entre pouvoir prédictif, stabilité et
interprétabilité. Une étude comparative confronte ces critères — mesurés par l'AUC-ROC, la
statistique de Kolmogorov--Smirnov, la précision, le rappel et la stabilité temporelle — afin
de justifier la sélection du modèle retenu pour la banque. L'écosystème \texttt{scikit-learn}
\cite{pedregosa2011} fournit un cadre unifié pour l'entraînement, la validation croisée et
l'évaluation de ces modèles.

\section{Explainable AI (XAI)}

L'explicabilité est une exigence réglementaire et éthique majeure dans le contexte bancaire,
particulièrement lorsque des modèles non linéaires sont mobilisés pour des décisions
individuelles~\cite{molnar2022}.

\subsection{Nécessité de l'explicabilité dans la finance réglementée}

Toute décision de crédit doit être justifiable et auditable. L'explicabilité répond à une
double exigence : réglementaire, pour satisfaire les régulateurs et auditeurs, et éthique,
pour garantir l'absence de discrimination et la confiance des clients. On distingue les
modèles intrinsèquement interprétables (régression logistique, arbres peu profonds) des
méthodes d'explicabilité \textit{post-hoc}, appliquées après coup à un modèle complexe.

\subsection{Méthodes post-hoc : LIME et SHAP}

Parmi les méthodes post-hoc, LIME~\cite{ribeiro2016} approxime localement le modèle par un
modèle linéaire interprétable au voisinage de l'instance à expliquer. SHAP
(\textit{SHapley Additive exPlanations})~\cite{lundberg2017}, fondée sur la théorie des jeux
coopératifs, attribue à chaque variable une contribution à la prédiction correspondant à sa
valeur de Shapley. Contrairement à LIME, SHAP offre des garanties théoriques d'additivité, de
cohérence et de précision locale, ce qui en fait l'outil central de notre couche
d'explicabilité~\cite{shapdoc}.

\subsubsection{Fondements en théorie des jeux : la valeur de Shapley}

L'idée sous-jacente à SHAP provient de la théorie des jeux coopératifs, et plus précisément de
la valeur de Shapley, proposée par Lloyd Shapley en 1953 pour répartir équitablement un gain
commun entre les membres d'une coalition, en fonction de la contribution individuelle de
chacun. SHAP transpose ce cadre à l'explication d'une prédiction : les « joueurs » sont les
variables décrivant un dossier de crédit, et le « gain » à répartir est l'écart entre la
prédiction moyenne du modèle sur l'ensemble du portefeuille (la valeur de référence, ou
\textit{base value}) et la prédiction obtenue pour ce dossier précis.

Concrètement, la contribution attribuée à chaque variable correspond à sa contribution
marginale moyenne : on imagine que les variables sont révélées au modèle une à une, dans un
ordre donné, et l'on mesure de combien la prédiction se déplace lorsque la variable considérée
est ajoutée à celles déjà connues. Comme ce déplacement dépend de l'ordre dans lequel les
variables sont révélées, la valeur de Shapley moyenne cette contribution sur l'ensemble des
ordres possibles, ce qui garantit un partage équitable, indépendant de tout choix arbitraire
d'ordonnancement.

Cette origine en théorie des jeux confère à SHAP des propriétés d'équité que ne possèdent pas
les heuristiques d'attribution \textit{ad hoc} : deux variables ayant un effet strictement
identique sur la prédiction reçoivent la même contribution (symétrie), une variable sans aucun
effet reçoit une contribution nulle (variable muette), et la somme des contributions
individuelles reconstitue exactement l'écart entre la prédiction moyenne et la prédiction du
dossier étudié (efficacité, ou additivité). C'est cette dernière propriété qui permet, dans
notre plateforme, de présenter côte à côte la valeur de référence et la PD finale, reliées par
un diagramme en cascade (\textit{waterfall}) des contributions de chaque variable.

\subsubsection{Unité des contributions SHAP : une précaution d'interprétation}

Un point souvent négligé lors de l'interprétation des valeurs SHAP concerne l'unité dans
laquelle elles s'expriment — celle-ci correspond à la sortie \emph{brute} du modèle, et non
nécessairement à une probabilité. Pour un modèle d'ensemble d'arbres entraîné en classification
(gradient boosting notamment), \texttt{shap.TreeExplainer} restitue par défaut des contributions
exprimées en \textit{log-odds} (le score brut avant passage par la fonction logistique), et non
en points de probabilité de défaut~\cite{shapdoc}. Une contribution de 0,8 signifie alors que la
variable décale le score en log-odds de 0,8 — un déplacement qui ne se traduit pas linéairement
en variation de PD, en particulier loin des extrémités de l'échelle de probabilité. Pour obtenir
des contributions directement interprétables par un comité de crédit en points de PD, il est
nécessaire d'invoquer l'explainer avec le paramètre \texttt{model\_output="probability"}, au prix
toutefois d'une additivité SHAP qui n'est plus alors garantie de façon exacte, mais seulement
approchée. Cette distinction d'unité doit être explicitée dans toute restitution des résultats
SHAP à un public métier, sous peine d'une lecture erronée de l'ampleur réelle de l'effet d'une
variable sur le risque.

\subsection{Explicabilité globale versus locale}

L'explicabilité globale décrit le comportement d'ensemble du modèle sur le portefeuille
(importance moyenne des variables, \textit{summary plots}), tandis que l'explicabilité locale
justifie une décision individuelle (\textit{force} et \textit{waterfall plots} par dossier). La
complémentarité de ces deux niveaux, offerte nativement par les valeurs SHAP, répond aux
besoins respectifs des gestionnaires de risque et des analystes.

\section{Agents IA et systèmes d'orchestration}

L'architecture d'agent IA permet d'orchestrer des composants prédictifs et explicatifs au
sein d'un pipeline décisionnel cohérent et traçable.

\subsection{Architecture et principes des agents IA}

Un agent IA orchestre une séquence de composants — perception, raisonnement, décision et
action — au sein d'un pipeline cohérent. Dans notre cas, l'agent coordonne l'exécution du
modèle, la génération des explications et l'application des règles de décision, en sélectionnant
dynamiquement les traitements à invoquer selon la requête de l'analyste.

\subsection{Seuils de confiance et mécanismes d'escalade humaine}

L'agent applique des seuils de confiance qui déterminent si une décision peut être proposée
automatiquement ou doit être escaladée vers un décideur humain. Ce mécanisme
\textit{human-in-the-loop} garantit que les cas incertains ou à fort enjeu sont soumis au
comité de crédit, l'IA restant un outil d'assistance et non de substitution.

\subsection{Traçabilité, logging et gouvernance}

Chaque décision et chaque résultat intermédiaire sont journalisés afin d'assurer une
traçabilité complète, socle de la gouvernance et de l'auditabilité du système.

\section{Grands modèles de langage (LLM) pour la finance}

Les LLM apportent une couche d'interaction et de raisonnement en langage naturel adaptée
aux besoins des analystes crédit.

\subsection{Fondements et capacités de raisonnement des LLM}

Les grands modèles de langage reposent sur l'architecture Transformer et le mécanisme
d'attention~\cite{vaswani2017}, et acquièrent, par pré-entraînement à grande échelle, une
capacité de raisonnement contextuel et de génération de texte cohérent à partir d'instructions
en langage naturel (\textit{few-shot learning})~\cite{brown2020}. Ils permettent de transformer
des sorties techniques (scores, contributions SHAP) en explications accessibles.

\subsection{NLP et LLM appliqués à l'analyse du risque}

Le traitement du langage naturel ouvre des cas d'usage variés en finance : synthèse de
dossiers, extraction d'information et interaction conversationnelle. Des modèles spécialisés,
tels que BloombergGPT~\cite{wu2023bloomberggpt}, illustrent l'intérêt d'adapter les LLM au
domaine financier. Dans notre plateforme, le LLM sert principalement à expliquer les sorties du
modèle et à répondre aux questions des analystes, en s'appuyant sur une inférence locale via
Ollama~\cite{ollama} afin de préserver la souveraineté des données.

\subsection{Intégration LLM dans des pipelines décisionnels}

L'intégration d'un LLM exige une séparation stricte des rôles : le modèle prédictif estime
le risque, tandis que le LLM se limite à expliquer et à reformuler. Ce principe garantit que
le LLM n'introduit aucune prédiction de risque non contrôlée et qu'aucune valeur numérique
n'est inventée ou altérée par la couche de langage.

\section{Tests de stress et simulation Monte Carlo}

L'estimation de la probabilité de défaut décrite précédemment repose sur des conditions
économiques « centrales », c'est-à-dire les plus probables. Or la gestion prudentielle du
risque exige également d'évaluer le comportement du portefeuille lorsque l'environnement se
dégrade fortement. C'est précisément l'objet des tests de stress, qui complètent la mesure du
risque attendu par une mesure du risque extrême.

\subsection{Principes et finalité des tests de résistance bancaire}

Les tests de résistance (\textit{stress tests}) consistent à soumettre un portefeuille à des
conditions économiques défavorables mais plausibles, afin de mesurer sa capacité à absorber
des pertes sans compromettre la solvabilité de l'établissement~\cite{bcbs2018}. Apparus comme un
instrument central de la supervision bancaire à la suite de la crise financière de 2008, ils ont
été institutionnalisés par les accords de Bâle~II et Bâle~III et sont aujourd'hui régulièrement
conduits par les autorités de régulation pour vérifier l'adéquation des fonds propres.

Sur le plan méthodologique, on distingue habituellement deux grandes familles d'approches.
L'analyse de sensibilité fait varier un seul facteur de risque à la fois (par
exemple une hausse du taux de défaut) pour isoler son impact marginal. L'analyse par
scénarios, plus réaliste, fait évoluer simultanément un ensemble cohérent de variables
macro-économiques — croissance du PIB, taux de chômage, taux d'intérêt, inflation — selon une
narration économique donnée. Ces scénarios peuvent être historiques, en rejouant un
épisode de crise passé, ou hypothétiques, en construisant une situation adverse
prospective. La construction de scénarios sévères mais crédibles constitue ainsi le c\oe ur
de la démarche.

Dans le contexte spécifique du risque de crédit PME, les tests de stress présentent un intérêt
particulier : la forte sensibilité de ce segment aux cycles macro-économiques et aux chocs
sectoriels rend l'évaluation en conditions adverses indispensable à une appréciation réaliste
des pertes potentielles.

\subsection{Simulation stochastique Monte Carlo}

Les approches par scénarios discrets, bien qu'utiles, ne couvrent qu'un nombre limité de
trajectoires économiques et ne permettent pas de quantifier la probabilité d'occurrence des
pertes. La simulation de Monte Carlo lève cette limite en adoptant une approche probabiliste :
plutôt que de tester quelques scénarios déterminés, elle explore un très grand nombre de
trajectoires aléatoires afin de reconstituer la distribution complète des pertes possibles du
portefeuille~\cite{glasserman2004}.

\subsubsection{Principe de la méthode}

Issue des travaux menés à Los Alamos dans les années 1940 et nommée en référence aux jeux de
hasard~\cite{metropolis1949}, la méthode de Monte Carlo repose sur la loi des grands nombres :
en répétant un très grand nombre de tirages aléatoires d'un phénomène stochastique, la moyenne
empirique des résultats converge vers l'espérance théorique recherchée. Appliquée au risque de
crédit, elle consiste à modéliser les variables d'entrée incertaines (facteurs
macro-économiques, paramètres de risque) par des lois de probabilité, puis à générer un grand
nombre de réalisations aléatoires de ces variables. Pour chaque réalisation, on recalcule les
pertes du portefeuille, ce qui produit au final une distribution empirique des pertes plutôt
qu'une valeur ponctuelle.

Formellement, soit $Y = g(X)$ la grandeur d'intérêt — perte de portefeuille, probabilité de
défaut simulée, etc. — où $X$ est un vecteur aléatoire de facteurs de risque de loi $\mathbb{P}_X$
et $g$ une fonction de valorisation potentiellement non linéaire et sans forme close (dans notre
cas, le pipeline complet ingénierie de variables~$\to$~préprocesseur~$\to$~modèle). À partir de
$n$ tirages indépendants et identiquement distribués $X^{(1)}, \dots, X^{(n)} \sim \mathbb{P}_X$,
l'estimateur de Monte Carlo de l'espérance $\mathbb{E}[Y]$ s'écrit :

\begin{equation}
  \hat{Y}_n = \frac{1}{n} \sum_{i=1}^{n} g\big(X^{(i)}\big).
  \label{eq:mc-estimator}
\end{equation}

Par la loi forte des grands nombres, $\hat{Y}_n \to \mathbb{E}[Y]$ presque sûrement lorsque
$n \to \infty$. Le théorème central limite précise en outre la vitesse de convergence : si
$\sigma^2 = \mathrm{Var}(Y) < \infty$, alors

\begin{equation}
  \sqrt{n}\,\big(\hat{Y}_n - \mathbb{E}[Y]\big) \;\xrightarrow{\ d\ }\; \mathcal{N}\big(0,\, \sigma^2\big),
  \label{eq:mc-clt}
\end{equation}

de sorte que l'erreur d'estimation décroît en $\mathcal{O}(\sigma / \sqrt{n})$, \emph{indépendamment
de la dimension} $d$ du vecteur de facteurs $X$. C'est précisément cette propriété qui rend Monte
Carlo supérieur aux méthodes d'intégration numérique déterministes (grilles, quadrature), dont
l'erreur se dégrade exponentiellement avec $d$ — le classique \textit{fléau de la
dimension}~\cite{glasserman2004}. C'est aussi ce qui justifie, en pratique, l'emploi de quelques
centaines à quelques milliers de tirages : doubler la précision de l'estimateur exige de
quadrupler le nombre de simulations, non de le doubler.

\subsubsection{Justification théorique en termes intuitifs}

Au-delà de sa formulation mathématique, l'intuition qui justifie la méthode de Monte Carlo peut
se résumer sans recourir au formalisme. Le principe repose sur la loi des grands nombres : si
l'on répète un très grand nombre de fois un tirage aléatoire représentatif d'un phénomène
incertain, la moyenne des résultats obtenus se rapproche de plus en plus de la valeur « réelle »
que l'on cherche à connaître, même lorsque celle-ci ne peut être calculée directement — parce
que le modèle sous-jacent (ici, la chaîne ingénierie de variables puis préprocesseur puis
modèle XGBoost) est trop complexe pour être analysé de façon analytique. Autrement dit, simuler
un grand nombre de dossiers de crédit « plausibles » et observer la distribution des
probabilités de défaut ainsi obtenues permet d'approcher, de façon empirique, la véritable
distribution du risque, sans qu'il soit nécessaire d'en connaître la formule.

Le théorème central limite apporte une seconde garantie, tout aussi intuitive : plus le nombre
de tirages augmente, plus l'estimation obtenue se stabilise autour de la valeur recherchée, et
l'ampleur de l'erreur restante diminue de façon prévisible et régulière à mesure que les
simulations s'accumulent. Un résultat remarquable, et déterminant pour son usage pratique, est
que cette vitesse de convergence ne dépend pas du nombre de variables prises en compte dans la
simulation : qu'un dossier de crédit comporte une dizaine ou plusieurs dizaines de variables
(chiffre d'affaires, trésorerie, taux d'intérêt, indicateurs macro-économiques, etc.), la
précision obtenue pour un même nombre de tirages reste comparable. C'est précisément ce qui
rend la méthode praticable pour des problèmes à de nombreuses variables, là où des approches
déterministes par quadrillage systématique de l'espace des possibles deviendraient rapidement
impraticables à mesure que le nombre de variables augmente.

Cette double garantie — convergence vers la bonne valeur, et vitesse de convergence
indépendante de la dimension du problème — est ce qui légitime, sur le plan théorique, l'usage
de quelques milliers de simulations seulement pour caractériser de façon fiable la distribution
du risque d'un dossier ou d'un portefeuille, sans qu'il soit nécessaire d'expliciter la formule
mathématique sous-jacente pour en comprendre la portée.

\subsubsection{Étapes de mise en \oe uvre}

Le déroulement type d'une simulation Monte Carlo appliquée au risque de crédit comprend la
spécification des distributions de probabilité des variables d'entrée et de leurs corrélations,
le tirage aléatoire d'un jeu de valeurs à chaque itération, le calcul de la perte du
portefeuille pour ce jeu de valeurs, la répétition de l'opération sur un grand nombre
d'itérations, puis l'agrégation des résultats en une distribution des pertes. La fiabilité des
estimations dépend du nombre de simulations : la précision s'améliore à mesure que les tirages
se multiplient, jusqu'à atteindre la convergence des indicateurs étudiés.

\subsubsection{Indicateurs de risque dérivés}

L'exploitation de la distribution simulée des pertes permet d'estimer plusieurs mesures de
risque clés. La perte attendue correspond à la moyenne de la distribution, tandis que
l'analyse du comportement en queue (\textit{tail risk}) renseigne sur les pertes
extrêmes peu probables mais critiques. La Value-at-Risk (VaR) mesure la perte maximale
susceptible d'être subie pour un niveau de confiance donné (par exemple 99~\%), et
l'Expected Shortfall (ou VaR conditionnelle) en complète l'analyse en quantifiant la
perte moyenne au-delà de ce seuil. Ces indicateurs offrent une vision quantifiée et
probabilisée du risque extrême du portefeuille, particulièrement adaptée à l'évaluation de
portefeuilles PME hétérogènes et sensibles aux chocs.

Formellement, pour un niveau de confiance $\alpha \in (0,1)$ (typiquement 95~\% ou 99~\%), la
Value-at-Risk est le quantile d'ordre $\alpha$ de la distribution de la perte (ou, dans notre cas,
de la probabilité de défaut simulée) $Y$ :

\begin{equation}
  \mathrm{VaR}_\alpha(Y) = \inf\big\{\, y \in \mathbb{R} \;:\; \mathbb{P}(Y \le y) \ge \alpha \,\big\}.
  \label{eq:var}
\end{equation}

À partir de l'échantillon simulé $\{Y_1, \dots, Y_n\}$, cette quantité s'estime naturellement par
le quantile empirique d'ordre $\alpha$ de l'échantillon trié. La VaR présente cependant une
limite bien connue : elle ne renseigne en rien sur l'ampleur des pertes \emph{au-delà} du seuil
$\alpha$. L'Expected Shortfall (ES) corrige ce défaut en moyennant les réalisations qui
dépassent effectivement ce seuil :

\begin{equation}
  \mathrm{ES}_\alpha(Y) = \mathbb{E}\big[\, Y \;\big|\; Y \ge \mathrm{VaR}_\alpha(Y) \,\big]
    \;\approx\; \frac{1}{|T_\alpha|} \sum_{i \in T_\alpha} Y_i,
    \qquad T_\alpha = \big\{\, i \;:\; Y_i \ge \mathrm{VaR}_\alpha(Y) \,\big\}.
  \label{eq:es}
\end{equation}

Contrairement à la VaR, l'ES est une mesure de risque \emph{cohérente} : elle est notamment
sous-additive, ce qui garantit qu'elle ne sous-estime jamais le risque d'un portefeuille
diversifié — une propriété que la VaR ne vérifie pas en général~\cite{glasserman2004}. C'est
pourquoi les cadres prudentiels récents tendent à privilégier l'ES comme mesure de référence du
risque extrême, la VaR restant néanmoins largement utilisée pour sa simplicité
d'interprétation.

\subsubsection{Apports et limites}

La principale force de la simulation Monte Carlo réside dans sa flexibilité : elle s'accommode
de distributions non normales, de relations non linéaires et de structures de corrélation
complexes entre facteurs de risque, là où les approches analytiques classiques atteignent
leurs limites. Ses limites tiennent essentiellement à son coût de calcul, qui croît avec le
nombre d'itérations, et à sa forte dépendance à la qualité des hypothèses de distribution et de
corrélation retenues — selon le principe bien connu du \textit{garbage in, garbage out}, des
hypothèses mal calibrées produisent des résultats trompeurs.

\section*{Conclusion}

Cet état de l'art a posé le socle scientifique du projet, du risque de crédit PME aux méthodes
d'apprentissage automatique, d'explicabilité et de simulation stochastique. Le chapitre
suivant présente l'architecture technique et la conception de la plateforme.

% ════════════════════════════════════════════════════════════
%  CHAPITRE 3 – CONCEPTION ET ARCHITECTURE DE LA SOLUTION
% ════════════════════════════════════════════════════════════
\chapter{Conception et architecture de la solution}

\section*{Introduction}

Après avoir posé, dans les chapitres précédents, le cadre métier du financement
des PME et le socle scientifique du projet, ce troisième chapitre est consacré à
la conception et à l'architecture technique de la plateforme effectivement
réalisée. Il fait le pont entre les besoins exprimés et l'implémentation : il décrit
l'architecture logique multi-couches, la structure concrète du projet (organisation
du code), les choix technologiques retenus, la conception des données, l'architecture
de l'API et des services métier, l'orchestration de l'agent IA, le module
d'explicabilité, le moteur de simulation Monte~Carlo, la modélisation UML, ainsi
que la planification des sprints et des releases.

L'objectif de ce chapitre est double : justifier les décisions d'architecture au
regard des exigences fonctionnelles et non fonctionnelles (performance, modularité,
auditabilité, explicabilité), et fournir une vue suffisamment précise du système
pour en garantir la reproductibilité et la maintenabilité.

% ────────────────────────────────────────────────────────────
\section{Vision architecturale globale}

\subsection{Présentation générale de la plateforme}

La plateforme est conçue selon une architecture modulaire en couches, qui
sépare strictement les responsabilités : ingestion et validation des données,
prédiction de la probabilité de défaut (PD), explicabilité (SHAP), simulation de
scénarios et de stress (Monte~Carlo), orchestration décisionnelle et interaction en
langage naturel. Cette séparation des préoccupations
(\textit{separation of concerns}) répond directement aux exigences non
fonctionnelles de maintenabilité et d'évolutivité formulées au chapitre~1 : chaque
composant peut être testé, remplacé ou réentraîné indépendamment des autres.

Le système est structuré autour de trois grands ensembles :

\begin{itemize}
  \item un backend applicatif exposant une API REST (FastAPI)~\cite{fastapi}, qui
        orchestre la chaîne validation $\rightarrow$ prédiction $\rightarrow$
        explication $\rightarrow$ scénario $\rightarrow$ synthèse en langage naturel ;
  \item un noyau analytique (artefact de modèle sérialisé, moteur SHAP,
        moteur Monte~Carlo, ingénierie de variables) découplé de la couche HTTP ;
  \item une interface analyste (tableau de bord web statique) consommant
        l'API et restituant prédictions, explications et simulations sous forme
        graphique.
\end{itemize}

\subsection{Architecture logique multi-couches}

L'architecture s'organise en cinq couches complémentaires, chacune assurant
une fonction précise du pipeline décisionnel. La figure~\ref{fig:archi_couches}
en présente la vue d'ensemble.

\begin{figure}[H]
  \centering
  \begin{tikzpicture}[node distance=6mm]
    \node[layerbox] (l5) {\textbf{Couche interaction LLM \& interface analyste}\\
      Tableau de bord web · dialogue en langage naturel · synthèse narrative};
    \node[layerbox, below=of l5] (l4) {\textbf{Couche agent IA \& orchestration}\\
      Enchaînement des traitements · seuils de confiance · escalade humaine};
    \node[layerbox, below=of l4] (l3) {\textbf{Couche explicabilité \& traçabilité}\\
      SHAP (\texttt{TreeExplainer}) · contributions locales/globales · logging};
    \node[layerbox, below=of l3] (l2) {\textbf{Couche modélisation \& prédiction}\\
      Feature engineering · préprocesseur · \texttt{predict\_proba} · classe de risque};
    \node[layerbox, below=of l2] (l1) {\textbf{Couche collecte \& ingestion des données}\\
      Acquisition du dossier PME · validation des variables · formatage};
    % flèches bidirectionnelles entre couches
    \foreach \a/\b in {l5/l4, l4/l3, l3/l2, l2/l1}{
      \draw[arrow, {Stealth[length=2mm]}-{Stealth[length=2mm]}] (\a) -- (\b);
    }
    % annotations de flux (centrées : n'altèrent pas l'axe de la figure)
    \node[above=3mm of l5, font=\scriptsize\itshape, text=DarkGray] {Analyste crédit};
    \node[below=3mm of l1, font=\scriptsize\itshape, text=DarkGray] {Données du dossier PME};
  \end{tikzpicture}
  \caption{Architecture logique multi-couches de la plateforme}
  \label{fig:archi_couches}
\end{figure}

\subsubsection{Couche collecte et ingestion des données}

Cette couche centralise l'acquisition, la validation et le formatage des
données du dossier de crédit PME en provenance des systèmes de la banque. Le service
de validation (\texttt{validation\_service}) vérifie la présence des variables
attendues par le modèle, signale les champs manquants (qui seront imputés) et les
champs non utilisés, sans jamais bloquer la chaîne : la robustesse prime, les
manques étant gérés en aval par imputation.

\subsubsection{Couche modélisation et prédiction}

Elle héberge l'artefact entraîné (préprocesseur + modèle) chargé et mis en cache au
démarrage, et expose l'estimation de la PD pour chaque dossier soumis. C'est le
c\oe ur prédictif : application de l'ingénierie de variables, transformation par le
préprocesseur, puis \texttt{predict\_proba} et classement du risque
(\textit{low / medium / high}).

\subsubsection{Couche explicabilité et traçabilité}

Elle calcule les contributions SHAP (\texttt{TreeExplainer}) sur la matrice
de variables transformées et restitue, pour chaque prédiction, les facteurs qui
augmentent ou diminuent le risque. Chaque explication est structurée (valeur de base,
contributions ordonnées, direction) pour être auditable.

\subsubsection{Couche agent IA et orchestration}

Elle coordonne l'enchaînement des traitements (validation, prédiction, explication,
scénario), applique les seuils de confiance et gère la composition des résultats en
une réponse unique. C'est le niveau qui matérialise la logique d'\textit{agent} :
décider quels outils invoquer en fonction de la requête de l'analyste.

\subsubsection{Couche interaction LLM et interface analyste}

Elle expose une interface conversationnelle permettant aux analystes d'interroger le
système en langage naturel. Un grand modèle de langage (LLM) reçoit uniquement
les sorties calculées par les couches précédentes et produit une synthèse narrative,
sans jamais recalculer ni inventer de valeur numérique.

\subsection{Flux de données et interactions entre composants}

Les composants communiquent via des interfaces standardisées (modèles Pydantic
typés). Le flux nominal, depuis la soumission d'un dossier jusqu'à la restitution
d'une explication narrative, garantit la cohérence et la traçabilité de bout en bout.
La figure~\ref{fig:flux_donnees} en donne la représentation.

\begin{figure}[H]
  \centering
  \begin{tikzpicture}[node distance=9mm and 8mm]
    % Rangée 1 (gauche vers droite) : ingestion et préparation
    \node[fluxio]                    (in)   {Dossier\\PME};
    \node[fluxbox, right=of in]      (val)  {Validation\\\texttt{\tiny Pydantic}};
    \node[fluxbox, right=of val]     (fe)   {Ingénierie\\de variables};
    \node[fluxbox, right=of fe]      (prep) {Préprocesseur\\\texttt{\tiny joblib}};

    % Rangée 2 (droite vers gauche) : prédiction et explication
    \node[fluxbox, below=of prep]    (pred) {Modèle XGBoost\\\texttt{\tiny predict\_proba}};
    \node[fluxbox, left=of pred]     (shap) {Explication SHAP\\\texttt{\tiny TreeExplainer}};
    \node[fluxbox, left=of shap]     (scen) {Scénarios\\\textit{\tiny what-if} / Monte Carlo};

    % Rangée 3 (gauche vers droite) : décision et restitution
    \node[fluxbox, below=of scen]    (agent){Agent IA\\seuils \& escalade};
    \node[fluxbox, right=of agent]   (llm)  {Synthèse LLM\\\texttt{\tiny Ollama}};
    \node[fluxio,  right=of llm]     (out)  {Décision\\explicable};

    % Enchaînement en serpentin
    \draw[fluxarrow] (in)   -- (val);
    \draw[fluxarrow] (val)  -- (fe);
    \draw[fluxarrow] (fe)   -- (prep);
    \draw[fluxarrow] (prep) -- node[fluxlab, right]{PD} (pred);
    \draw[fluxarrow] (pred) -- (shap);
    \draw[fluxarrow] (shap) -- (scen);
    \draw[fluxarrow] (scen) -- node[fluxlab, left]{$\Delta$PD} (agent);
    \draw[fluxarrow] (agent)-- (llm);
    \draw[fluxarrow] (llm)  -- (out);

    % Traçabilité : chaque étape est journalisée
    \node[fluxlab, below=5mm of out, align=center] (jrn)
      {journalisation\\\& audit};
    \draw[fluxarrow, dashed] (out) -- (jrn);
  \end{tikzpicture}
  \caption{Flux de données de bout en bout du pipeline décisionnel}
  \label{fig:flux_donnees}
\end{figure}

% ────────────────────────────────────────────────────────────
\section{Structure du projet}

L'organisation du dépôt traduit directement l'architecture en couches : le backend
applicatif est isolé de l'interface, des notebooks de modélisation, des jeux de
données et des scripts d'entraînement. Le listing~\ref{lst:arbo} présente
l'arborescence simplifiée du projet, réduite à ses cinq répertoires essentiels.

\begin{codeblock}
\begin{lstlisting}[caption={Arborescence simplifiée du projet}, label={lst:arbo},
                   basicstyle=\ttfamily\scriptsize, columns=fixed, numbers=none]
BFPME/
+-- backend/                   # Application FastAPI (API REST + services)
|   +-- app/
|   |   +-- main.py            # App FastAPI, montage des routers + frontend
|   |   +-- config.py          # Configuration Pydantic (variables .env)
|   |   +-- schemas.py         # Modeles de requete/reponse (Pydantic)
|   |   +-- routers/           # Un fichier par groupe d'endpoints
|   |   |   +-- prediction.py  explain.py  scenario.py
|   |   |   +-- monte_carlo.py chat.py     metadata.py
|   |   +-- services/          # Logique metier decouplee du HTTP
|   |       +-- artifacts.py            # Chargement/cache de l'artefact joblib
|   |       +-- model_service.py        # Feature engineering + predict_proba
|   |       +-- shap_service.py         # TreeExplainer SHAP
|   |       +-- scenario_service.py     # PD baseline vs. modifiee (delta_pd)
|   |       +-- monte_carlo_service.py  # Stress test stochastique vectorise
|   |       +-- llm_service.py          # Appel async httpx vers Ollama
|   |       +-- feature_engineering.py  # Ratios derives a l'inference
|   +-- requirements.txt
+-- frontend/                  # Tableau de bord statique (HTML / CSS / JS)
|   +-- index.html  app.js  styles.css
+-- notebook/                  # Notebooks de modelisation (Options A / B / C)
+-- data/                      # Datasets + artefacts modeles (.joblib, .json)
+-- scripts/                   # Entrainement et generation de donnees
    +-- train_option_c_xgboost.py
\end{lstlisting}
\end{codeblock}

Ce découpage applique le principe de séparation entre la couche HTTP et la
logique métier : les fichiers de \texttt{routers/} ne contiennent que la déclaration
des points d'entrée et la validation des entrées/sorties, tandis que toute la logique
analytique réside dans \texttt{services/}. Un service peut ainsi être réutilisé par
plusieurs endpoints (par exemple \texttt{model\_service.predict} est appelé à la fois
par \texttt{/predict}, \texttt{/scenario} et \texttt{/chat/analyze}).

% ────────────────────────────────────────────────────────────
\section{Choix technologiques et environnement de développement}

Le tableau~\ref{tab:techno} récapitule la pile technologique retenue et la
justification de chaque choix au regard des besoins du projet.

\begin{table}[H]
\centering
\caption{Pile technologique de la plateforme}
\label{tab:techno}
\renewcommand{\arraystretch}{1.3}
\begin{tabularx}{\textwidth}{@{}l l Y@{}}
\toprule
\textbf{Domaine} & \textbf{Technologie} & \textbf{Justification} \\
\midrule
Langage         & Python 3            & Écosystème data science / ML mature \\
API REST        & FastAPI + Uvicorn   & Asynchrone, typage natif, OpenAPI/Swagger auto \\
Validation      & Pydantic v2         & Schémas typés, configuration par variables d'env. \\
ML --- arbres   & XGBoost, scikit-learn (RandomForest) & Modèles ensemblistes performants \\
Prétraitement   & scikit-learn        & \texttt{ColumnTransformer}, encodeurs, scalers \\
Explicabilité   & SHAP (\texttt{TreeExplainer}) & Explications cohérentes globales et locales \\
Sérialisation   & joblib              & Persistance préprocesseur + modèle \\
LLM             & Ollama (API compatible OpenAI) & Inférence locale, souveraineté des données \\
Client HTTP     & httpx (async)       & Appels non bloquants vers le LLM \\
Frontend        & HTML/CSS/JS, Chart.js & Tableau de bord sans étape de build \\
\bottomrule
\end{tabularx}
\end{table}

\subsection{Configuration par variables d'environnement}

Toute la configuration (chemins des artefacts, seuil de décision, activation et
paramétrage du LLM, hôte/port) est centralisée dans \texttt{config.py} via
\texttt{BaseSettings} de Pydantic~\cite{pydantic}, alimentée par un fichier \texttt{.env}. Ce
mécanisme découple le code de l'infrastructure : il permet de basculer de modèle
(RandomForest $\leftrightarrow$ XGBoost), d'activer ou non le LLM, ou de durcir le
seuil de décision sans modifier le code.

\begin{codeblock}
\begin{lstlisting}[caption={Extrait de configuration (.env)}, label={lst:env}]
MODEL_ARTIFACT_PATH=d:/BFPME/data/option_c_xgboost_model.joblib
MODEL_METRICS_PATH=d:/BFPME/data/option_c_xgboost_metrics.json
DEFAULT_THRESHOLD=0.5            # PD >= seuil => defaut predit
LLM_ENABLED=true
LLM_BASE_URL=http://127.0.0.1:11434/v1   # Ollama
LLM_MODEL=qwen2.5:7b-instruct
LLM_MAX_TOKENS=500
\end{lstlisting}
\end{codeblock}

% ────────────────────────────────────────────────────────────
\section{Conception des données}

\subsection{Sources et nature des données bancaires}

Les données proviennent des dossiers de crédit PME : informations signalétiques de
l'entreprise (ancienneté, secteur, région, structure juridique, classification PME,
effectif), états financiers (chiffre d'affaires, total actif/passif, capitaux
propres, dettes court et long terme, fonds de roulement, BFR, trésorerie), variables
de crédit (montant du prêt, durée, type de crédit, objectif, type de taux, valeur des
garanties), variables comportementales (retards de paiement, impayés) et variables
de scoring interne BFPME. La variable cible est \texttt{default\_flag} (défaut /
non-défaut).

Conformément au cadre du projet, le jeu de données exploité
(\texttt{final\_dataset.csv}) est synthétique mais métier-cohérent : il est
généré à partir d'un \textit{driver} latent de risque dérivé du score BFPME, de la
classe de risque et du type de projet (création / extension), de sorte que les
variables financières, comportementales et de garantie restent corrélées de façon
réaliste. Il sert de support de prototypage, de comparaison de modèles et
d'ingénierie de variables.

\subsection{Modèle de données et entités principales}

Le modèle de données structure les entités centrales du domaine du crédit PME et
leurs relations. La figure~\ref{fig:classes_donnees} (diagramme de classes, non
modifié) en formalise la structure ; les entités principales sont :

\begin{itemize}
  \item Entreprise / PME : caractéristiques signalétiques et financières du
        demandeur ;
  \item Dossier de crédit : montant, durée, type, objectif, garanties ;
  \item Décision : PD estimée, classe de risque, seuil appliqué, explication
        associée ;
  \item Scénario de stress : chocs appliqués, distribution de PD simulée et
        indicateurs de risque dérivés.
\end{itemize}

\subsection{Pipeline de prétraitement (artefact joblib)}

Le préprocesseur et le modèle sont entraînés conjointement puis sérialisés dans un
unique artefact joblib contenant : le \texttt{preprocessor}
(\texttt{ColumnTransformer} ajusté), le \texttt{model} entraîné et la liste
\texttt{allowed\_columns} des colonnes attendues. Le préprocesseur applique, selon le
type de variable, la chaîne suivante :

\begin{enumerate}
  \item variables numériques sensibles aux valeurs extrêmes : imputation par
        la médiane $\rightarrow$ winsorisation (écrêtage aux centiles 1 et 99)
        $\rightarrow$ standardisation ;
  \item variables numériques simples : imputation médiane $\rightarrow$
        standardisation ;
  \item variables ordinales (ex. durée de crédit) : imputation par mode
        $\rightarrow$ \texttt{OrdinalEncoder} avec ordre métier explicite ;
  \item variables nominales (secteur, région, structure juridique, type de
        crédit, scénario économique) : imputation par mode $\rightarrow$
        \texttt{OneHotEncoder} (\texttt{handle\_unknown=ignore}).
\end{enumerate}

Ce découpage par type de variable, encapsulé dans un \texttt{ColumnTransformer},
porte la dimensionnalité de 52~variables brutes à 80~variables transformées.
La figure~\ref{fig:pipeline_prepro} en illustre le déroulé.

\begin{figure}[H]
  \centering
  \resizebox{\textwidth}{!}{%
  \begin{tikzpicture}[node distance=5mm and 9mm]
    % entrée
    \node[flowbox, fill=EspritRed!12] (raw) {Variables\\brutes\\(52 var.)};
    % quatre branches par type de variable
    \node[flowbox, right=of raw, yshift=2.55cm] (num1) {Num. extrêmes};
    \node[flowbox, right=of raw, yshift=0.85cm] (num2) {Num. simples};
    \node[flowbox, right=of raw, yshift=-0.85cm] (ord)  {Ordinales};
    \node[flowbox, right=of raw, yshift=-2.55cm] (nom)  {Nominales};
    % traitements
    \node[flowbox, right=of num1] (t1) {Impute méd. $\rightarrow$\\Winsor. $\rightarrow$\\Standard.};
    \node[flowbox, right=of num2] (t2) {Impute méd. $\rightarrow$\\Standard.};
    \node[flowbox, right=of ord]  (t3) {Impute mode $\rightarrow$\\OrdinalEnc.};
    \node[flowbox, right=of nom]  (t4) {Impute mode $\rightarrow$\\OneHotEnc.};
    % concaténation
    \node[flowbox, right=16mm of t2, yshift=-0.85cm, fill=EspritRed!12,
          minimum height=3.4cm] (concat) {Concaténation\\(\texttt{Column-}\\\texttt{Transformer})};
    \node[flowbox, right=of concat] (out) {Matrice\\transformée\\(80 var.)};
    % flèches
    \foreach \b in {num1,num2,ord,nom}{\draw[arrow] (raw) -- (\b);}
    \draw[arrow] (num1) -- (t1); \draw[arrow] (num2) -- (t2);
    \draw[arrow] (ord)  -- (t3); \draw[arrow] (nom)  -- (t4);
    \foreach \t in {t1,t2,t3,t4}{\draw[arrow] (\t) -- (concat);}
    \draw[arrow] (concat) -- (out);
  \end{tikzpicture}}
  \caption{Pipeline de prétraitement sérialisé dans l'artefact joblib}
  \label{fig:pipeline_prepro}
\end{figure}

\subsection{Ingénierie de variables à l'inférence}

Un choix de conception déterminant consiste à appliquer l'ingénierie de
variables avant le préprocesseur, à l'inférence, et non à la figer dans l'artefact.
Le service \texttt{feature\_engineering} dérive à la volée des ratios métier
discriminants à partir des champs bruts du dossier. Les principaux ratios sont
définis comme suit :

\begin{align}
  \text{loan\_to\_revenue}
    &= \frac{\text{montant\_pret}}{\text{chiffre\_affaires\_annuel}}, \\[4pt]
  \text{cash\_to\_debt}
    &= \frac{\text{tresorerie\_disponible}}{\text{dette\_totale}}, \\[4pt]
  \text{garantie\_to\_loan}
    &= \frac{\text{valeur\_garanties}}{\text{montant\_pret}}, \\[4pt]
  \text{payment\_stress\_index}
    &= \mathrm{clip}\!\left(\tfrac{\text{retards\_jours}}{90},\,0,\,3\right)
     + \mathrm{clip}\!\left(\text{impayes\_to\_loan},\,0,\,3\right).
\end{align}

D'autres ratios sont calculés selon la même logique (\texttt{debt\_total},
\texttt{impayes\_to\_loan}, \texttt{fdr\_bfr\_pressure},
\texttt{financement\_externe\_intensity}). Cette approche garantit que la
même logique de dérivation est appliquée à l'entraînement, à la prédiction
unitaire, au scénario et à la simulation Monte~Carlo, évitant toute dérive entre
apprentissage et production.

% ────────────────────────────────────────────────────────────
\section{Conception de l'API et des services}

\subsection{Points d'entrée REST}

L'API expose un ensemble cohérent de points d'entrée, chacun associé à un routeur
dédié. Le tableau~\ref{tab:endpoints} en présente la cartographie.

\begin{table}[H]
\centering
\caption{Cartographie des points d'entrée de l'API}
\label{tab:endpoints}
\renewcommand{\arraystretch}{1.3}
\begin{tabularx}{\textwidth}{@{}l l Y@{}}
\toprule
\textbf{Méthode} & \textbf{Endpoint} & \textbf{Rôle} \\
\midrule
GET  & \texttt{/health}              & Sonde de disponibilité du service \\
GET  & \texttt{/metadata/model}      & Métadonnées et métriques du modèle \\
GET  & \texttt{/metadata/eda}        & Résumé exploratoire du dataset (graphes) \\
POST & \texttt{/metadata/reload-model} & Rechargement de l'artefact sans redémarrage \\
POST & \texttt{/validate}            & Contrôle de cohérence du dossier \\
POST & \texttt{/predict}             & Prédiction de la PD et de la classe de risque \\
POST & \texttt{/explain}             & Contributions SHAP (top-$k$) \\
POST & \texttt{/scenario}            & Simulation what-if (PD baseline vs. modifiée) \\
POST & \texttt{/monte-carlo}         & Test de stress stochastique \\
POST & \texttt{/chat/analyze}        & Orchestration complète + synthèse LLM \\
\bottomrule
\end{tabularx}
\end{table}

\subsection{Cycle de vie de l'artefact}

L'artefact est chargé une seule fois au démarrage de l'application
(événement \texttt{startup}) et mis en cache en mémoire (\texttt{lru\_cache}). Le
point d'entrée \texttt{/metadata/reload-model} permet de vider ce cache et de
recharger le modèle à chaud, sans interruption de service, après un réentraînement.
Cette conception répond à l'exigence de performance (pas de rechargement disque à
chaque requête) tout en préservant la maintenabilité opérationnelle.

% ────────────────────────────────────────────────────────────
\section{Architecture de l'agent IA et orchestration}

\subsection{Pipeline d'orchestration}

Le c\oe ur de l'agent réside dans l'endpoint \texttt{/chat/analyze}, qui orchestre
les outils analytiques en fonction des indicateurs portés par la requête
(\texttt{run\_validation}, \texttt{run\_prediction}, \texttt{run\_explanation},
\texttt{scenario\_changes}). Il compose les résultats dans une structure unique
\texttt{tool\_results}, puis délègue la rédaction de la synthèse au LLM. La
figure~\ref{fig:seq_chat} (diagramme de séquence, non modifié) en détaille les
échanges.

\begin{figure}[H]
  \centering
  \includegraphics[width=0.95\textwidth]{sequence_chat_analyze}
  \caption{Diagramme de séquence de l'orchestration \texttt{/chat/analyze}}
  \label{fig:seq_chat}
\end{figure}

\begin{codeblock}
\begin{lstlisting}[caption={Logique d'orchestration (extrait)}, label={lst:orch},
                   language=Python]
if run_validation:  tool_results["validation"]  = validate_features(features)
if run_prediction:  tool_results["prediction"]  = predict(features, threshold)
if run_explanation: tool_results["explanation"] = explain_prediction(features, k)
if scenario_changes: tool_results["scenario"]   = run_scenario(features, changes)
# synthese narrative deleguee au LLM (si active), sinon sorties brutes
\end{lstlisting}
\end{codeblock}

\subsection{Seuils de confiance et classes de risque}

L'agent applique un seuil de décision configurable (\texttt{DEFAULT\_THRESHOLD},
0{,}5 par défaut) : la classe prédite vaut « défaut » dès que la PD dépasse ce seuil.
Indépendamment, une classe de risque qualitative est attribuée à des fins d'escalade :
\textit{low} (PD~$<$~0{,}4), \textit{medium} ($0{,}4 \le$~PD~$<$~0{,}7) et
\textit{high} (PD~$\ge$~0{,}7). Ce double mécanisme — seuil quantitatif et classe
qualitative — matérialise la logique \textit{human-in-the-loop} : les dossiers à
risque élevé sont signalés pour escalade vers le comité de crédit, l'IA assistant la
décision sans s'y substituer.

\subsection{Séparation stricte des rôles modèle / LLM}

Un principe de conception fondamental gouverne l'intégration du LLM : la
séparation stricte des responsabilités. Le modèle prédictif estime le risque ;
le LLM se limite à expliquer et reformuler en langage naturel à partir des
sorties calculées. Le \textit{prompt} système impose explicitement de n'utiliser que
les sorties fournies et de ne jamais inventer, arrondir ou estimer une valeur
numérique. L'appel est asynchrone (\texttt{httpx}), à température basse (0{,}1) pour
limiter la variabilité, et toute défaillance du LLM est interceptée sans
jamais invalider les résultats déjà calculés — l'analyste reçoit alors les sorties
brutes du modèle, garantissant la continuité de service.

% ────────────────────────────────────────────────────────────
\section{Module d'explicabilité SHAP}

Le service d'explicabilité instancie un \texttt{TreeExplainer} (mis en cache) sur le
modèle entraîné et calcule les valeurs SHAP sur la matrice de variables
transformées. Pour chaque prédiction, il extrait les $k$ contributions de plus forte
amplitude, indique pour chacune sa direction (\textit{increase\_risk} /
\textit{decrease\_risk}) et restitue la valeur de base (\textit{base value}). Cette
restitution sert à la fois l'explicabilité locale (justification d'une
décision individuelle) et, par agrégation, l'explicabilité globale
(importance des variables sur le portefeuille), répondant aux besoins respectifs des
analystes et des gestionnaires de risque.

% ────────────────────────────────────────────────────────────
\section{Moteur de simulation Monte Carlo}

Le moteur de stress test génère, à partir d'un client de base et d'un ensemble de
chocs sur des variables choisies, $N$ scénarios perturbés, puis exécute le
modèle sur l'ensemble du lot en une seule passe vectorisée
(un \texttt{transform} + un \texttt{predict\_proba}), ce qui garantit la performance.
Les chocs sont appliqués aux variables \emph{brutes} (selon une loi normale ou
uniforme), puis les ratios d'ingénierie sont re-dérivés pour rester cohérents.

La distribution de PD obtenue est synthétisée par des indicateurs de risque
directement exploitables : moyenne et écart-type, percentiles, Value-at-Risk
à 95~\% (\textit{VaR}$_{95}$), Expected Shortfall (perte moyenne au-delà de
la VaR), probabilité de dépassement du seuil, répartition par classe de risque,
histogramme et scénarios extrêmes. Ces mesures offrent une lecture probabilisée de la
robustesse du portefeuille PME face aux chocs macro-économiques et sectoriels.

% ────────────────────────────────────────────────────────────
\section{Modélisation UML}

La conception du système a été formalisée à l'aide du langage UML. Les diagrammes
présentés ci-dessous ne sont pas modifiés et sont repris tels qu'élaborés
lors de la phase de conception.

\subsection{Diagramme de cas d'utilisation}

Le diagramme de cas d'utilisation (figure~\ref{fig:uc}) identifie les acteurs
principaux — l'analyste crédit, le gestionnaire de risque et le comité de crédit —
et leurs interactions avec la plateforme : prédire la PD, expliquer une décision,
simuler un scénario, lancer un test de stress et dialoguer en langage naturel.

\begin{figure}[H]
  \centering
  \includegraphics[width=0.9\textwidth]{diagramme_cas_utilisation}
  \caption{Diagramme de cas d'utilisation de la plateforme}
  \label{fig:uc}
\end{figure}

\subsection{Diagramme de classes}

Le diagramme de classes (figure~\ref{fig:classes_donnees}) formalise les entités du
domaine et les composants de service (préprocesseur, modèle, explicateur,
orchestrateur) ainsi que leurs relations.

\begin{figure}[H]
  \centering
  \includegraphics[width=0.95\textwidth]{diagramme_classes}
  \caption{Diagramme de classes du système}
  \label{fig:classes_donnees}
\end{figure}

\subsection{Diagramme de déploiement}

Le diagramme de déploiement (figure~\ref{fig:deploiement}) situe les n\oe uds
d'exécution : le serveur applicatif (FastAPI/Uvicorn), le serveur d'inférence LLM
local (Ollama) et le navigateur client hébergeant le tableau de bord.

\begin{figure}[H]
  \centering
  \includegraphics[width=0.9\textwidth]{diagramme_deploiement}
  \caption{Diagramme de déploiement de la plateforme}
  \label{fig:deploiement}
\end{figure}

% ────────────────────────────────────────────────────────────
\section{Planification des sprints et des releases}

Conformément au choix méthodologique SCRUM (chapitre~1), le développement a été
organisé en releases incrémentales, chacune livrant un sous-ensemble cohérent et
validable de la plateforme. Le tableau~\ref{tab:releases} synthétise cette
planification.

\begin{table}[H]
\centering
\caption{Planification des releases du projet}
\label{tab:releases}
\renewcommand{\arraystretch}{1.3}
\begin{tabularx}{\textwidth}{@{}l Y l@{}}
\toprule
\textbf{Release} & \textbf{Périmètre} & \textbf{Chapitre} \\
\midrule
Release~I   & Compréhension métier et ingénierie des données (cycle du crédit, collecte, nettoyage, \textit{feature engineering}) & Ch.~4 \\
Release~II  & Modélisation prédictive (régression logistique, arbre de décision, Random Forest, Gradient Boosting) et explicabilité SHAP & Ch.~5 \\
Release~III & Agent IA, orchestration décisionnelle et intégration LLM en langage naturel & Ch.~6 \\
Release~IV  & Tests de stress du portefeuille par simulation Monte~Carlo & Ch.~7 \\
\bottomrule
\end{tabularx}
\end{table}

Chaque release a été clôturée par une revue de sprint avec les parties prenantes de
MPSoft, suivie d'une rétrospective, garantissant un rythme de livraison régulier et
une adaptation continue aux retours métier.

\section*{Conclusion}

Ce chapitre a présenté la conception et l'architecture technique de la plateforme :
une architecture en cinq couches matérialisée par une structure de projet claire
séparant la couche HTTP de la logique métier ; une pile technologique (FastAPI,
scikit-learn/XGBoost, SHAP, Ollama) cohérente avec les exigences de performance,
d'explicabilité et de souveraineté des données ; une conception des données
articulant un préprocesseur sérialisé et une ingénierie de variables à l'inférence ;
et une orchestration d'agent IA assurant une séparation stricte entre estimation du
risque et restitution en langage naturel. La modélisation UML et la planification des
sprints complètent cette vue de conception. Les chapitres suivants détaillent la
réalisation de chacune des releases, en commençant par la compréhension métier et
l'ingénierie des données.

% ════════════════════════════════════════════════════════════
%  CHAPITRE 4 – RELEASE I
% ════════════════════════════════════════════════════════════
\chapter{Release~I : Compréhension du risque de crédit PME et ingénierie des données}

\section*{Introduction}

Cette première release pose les fondations du projet : la compréhension du métier du crédit
PME et la construction d'un jeu de données de qualité. En l'absence de données réelles
mobilisables pour des raisons de confidentialité et de protection des données, l'enjeu
central de cette release a été de générer un jeu de données synthétique suffisamment
cohérent, riche et réaliste pour servir de support fiable au prototypage, à la comparaison de
modèles et à l'ingénierie de variables. Les trois sprints qui composent cette release couvrent
respectivement la compréhension métier, la construction et l'exploration du jeu de données, puis
son prétraitement et son enrichissement.

\section{Sprint~1 : Compréhension métier et cycle de vie du crédit PME}

\subsection{Backlog du sprint}

Le backlog de ce sprint vise à formaliser le processus de crédit de la banque et à identifier
les risques à modéliser. Les \textit{user stories} retenues couvrent la cartographie du cycle
de vie d'un dossier PME, la distinction des catégories de risque et le recueil des besoins des
parties prenantes.

\subsection{Cycle de vie du crédit en banque}

Le cycle de vie d'un crédit PME s'articule en quatre grandes phases, de la demande initiale
jusqu'au recouvrement éventuel. La figure~\ref{fig:cycle_credit} en donne une représentation
synthétique.

\begin{figure}[H]
  \centering
  \resizebox{\textwidth}{!}{%
  \begin{tikzpicture}[node distance=10mm]
    \node[stagebox] (s1) {\textbf{1. Demande \&\\instruction}\\
      Dépôt du dossier PME · éligibilité · collecte des variables};
    \node[stagebox, right=of s1] (s2) {\textbf{2. Analyse \&\\décision}\\
      Évaluation · estimation de la PD · comité de crédit};
    \node[stagebox, right=of s2] (s3) {\textbf{3. Décaissement\\\& suivi}\\
      Octroi · suivi des remboursements · signaux comportementaux};
    \node[stagebox, right=of s3] (s4) {\textbf{4. Défaut \&\\recouvrement}\\
      Impayés persistants · \texttt{default\_flag} · recouvrement};
    \draw[arrow] (s1) -- (s2);
    \draw[arrow] (s2) -- (s3);
    \draw[arrow] (s3) -- (s4);
    % boucle de rétroaction des données vers la modélisation
    \draw[arrow, dashed] (s4.south) -- ++(0,-0.7) -| (s2.south)
      node[pos=0.25, below, font=\scriptsize, text=DarkGray]
      {Historique de défaut $\rightarrow$ réentraînement du modèle};
  \end{tikzpicture}}
  \caption{Cycle de vie du crédit PME et point d'insertion de la modélisation de la PD}
  \label{fig:cycle_credit}
\end{figure}

\subsubsection{Demande et instruction du dossier PME}

La PME soumet un dossier comprenant ses états financiers et son projet ; l'analyste vérifie
la complétude et l'éligibilité. C'est à cette étape que sont collectées les variables
signalétiques, financières et de garantie qui alimenteront ultérieurement le modèle.

\subsubsection{Analyse et décision du comité de crédit}

Le dossier est évalué puis soumis au comité de crédit qui statue sur l'octroi. La plateforme
vise à enrichir cette étape par une estimation de PD explicable, sans se substituer à la
décision du comité.

\subsubsection{Décaissement et suivi}

Après accord, le crédit est décaissé et fait l'objet d'un suivi régulier. Le comportement de
remboursement observé durant cette phase (retards, impayés, utilisation des lignes) constitue
une source majeure de signaux de risque.

\subsubsection{Défaut et recouvrement}

En cas d'impayé persistant, le dossier entre en phase de recouvrement. L'événement de défaut
matérialisé à ce stade correspond à la variable cible \texttt{default\_flag} du modèle.

\subsection{Identification des risques : éligibilité et PD}

Le sprint distingue le risque d'éligibilité, traité par des règles métier, de la probabilité
de défaut, objet de la modélisation. Cette distinction structure l'ensemble du projet :
l'éligibilité relève de filtres déterministes, tandis que la PD est un problème de
classification probabiliste supervisée.

\subsection{Parties prenantes et besoins}

Les besoins des analystes, gestionnaires et comités sont recueillis, notamment l'exigence
d'explicabilité et de traçabilité des décisions. Ces besoins orientent à la fois le choix des
variables (financières, sectorielles, comportementales et macro-économiques) et l'architecture
de la plateforme.

\subsection{Résultats du sprint}

Le sprint aboutit à une cartographie complète du processus et des risques, validée par les
parties prenantes, ainsi qu'à la spécification des familles de variables à reconstituer dans
le jeu de données synthétique.

\section{Sprint~2 : Collecte, génération et exploration des données bancaires}

\subsection{Backlog du sprint}

Ce sprint vise à constituer et explorer le jeu de données. En l'absence d'accès à des données
de production, il a consisté à concevoir un générateur de données synthétiques capable
de reproduire la structure et les corrélations métier d'un portefeuille de crédit PME, puis à
explorer le jeu obtenu.

\subsection{Sources et familles de données}

\subsubsection{Données clients et financières des PME}

Variables signalétiques (structure juridique, effectif, classification PME, ancienneté,
localisation) et états financiers (chiffre d'affaires, capitaux propres, dettes court et long
terme, FDR, BFR, trésorerie, EBITDA, marges, résultat net).

\subsubsection{Données des prêts et historique de remboursement}

Caractéristiques des crédits (montant, durée, type, objectif, taux, garanties) et comportement
de remboursement (DPD moyen et maximal, fréquence des retards, impayés, respect des échéances).

\subsubsection{Données comportementales et macro-économiques}

Indicateurs comportementaux bancaires (découverts fréquents, baisse des dépôts, utilisation des
lignes de crédit) et variables de contexte sectoriel et macro-économique (inflation, taux
directeur, croissance du PIB, chômage, risque sectoriel et pays).

\subsection{Génération du jeu de données synthétique}

Le jeu de données est construit en deux temps : un socle de scoring BFPME, puis un
enrichissement métier piloté par un facteur de risque latent. L'ensemble du processus
est rendu reproductible par une graine aléatoire fixe (\texttt{SEED~=~42}). La
figure~\ref{fig:pipeline_generation} en résume les étapes.

\begin{figure}[H]
  \centering
  \begin{tikzpicture}[node distance=5.5mm,
     stage/.style = {draw=DarkGray, thick, rounded corners=3pt, fill=LightGray,
                     text width=10cm, minimum height=0.95cm, align=center, font=\small},
     io/.style    = {draw=DarkGray, very thick, rounded corners=7pt, fill=white,
                     text width=8.5cm, minimum height=0.8cm, align=center,
                     font=\small\bfseries},
     down/.style  = {-{Stealth[length=2.5mm]}, thick, DarkGray}]

    \node[io] (in) {Tables de référence BFPME\\
      \texttt{\footnotesize RU\_Facteur · RU\_SousFacteur · RU\_CritereNotes}};
    \node[stage, below=of in] (s1)
      {\textbf{Étape 1 — Socle de scoring BFPME}\\
       scores, \texttt{risk\_class} et 62 sous-facteurs \texttt{SF\_*}};
    \node[stage, below=of s1] (s2)
      {\textbf{Étape 2 — Driver latent de risque}\\
       \texttt{latent\_risk} / \texttt{latent\_quality} + bruit gaussien};
    \node[stage, below=of s2] (s3)
      {\textbf{Étape 3 — Génération conditionnelle}\\
       ${\sim}90$ variables métier corrélées au profil latent};
    \node[stage, below=of s3] (s4)
      {\textbf{Étape 4 — Bruit, stochasticité \& écrêtage}\\
       tirages Bernoulli / multinomiaux, \texttt{clip} des domaines};
    \node[stage, below=of s4] (s5)
      {\textbf{Étape 5 — Logique de défaut synthétique}\\
       $\text{pd\_target}=\sigma(z)$,\; \texttt{default\_flag} $\sim$ Bernoulli};
    \node[io, below=of s5] (out)
      {\texttt{final\_dataset.csv}\; (cible \texttt{default\_flag})};

    \foreach \a/\b in {in/s1, s1/s2, s2/s3, s3/s4, s4/s5, s5/out}{
      \draw[down] (\a) -- (\b);
    }

    % Note latérale : reproductibilité
    \node[right=7mm of s3, text width=2.4cm, align=center,
          font=\scriptsize\itshape, text=DarkGray] (seed)
      {Processus reproductible\\ \texttt{SEED = 42}};
    \draw[dashed, DarkGray, -{Stealth[length=1.6mm]}] (seed) -- (s3.east);
  \end{tikzpicture}
  \caption{Pipeline de génération du jeu de données synthétique}
  \label{fig:pipeline_generation}
\end{figure}

\subsubsection{Étape 1 --- Socle de scoring BFPME}

Le point de départ (\texttt{bfpme\_synthetic\_dataset.csv}) est un ensemble de
1~000~dossiers de financement synthétiques, scorés selon la formule officielle
du moteur de scoring BFPME (reproduite depuis \texttt{ScoringService.cs}) à partir des tables
de référence (\texttt{RU\_Facteur}, \texttt{RU\_SousFacteur}, \texttt{RU\_CritereNotes}).
Ce socle fournit le code du modèle de scoring (\texttt{code\_modele} : création vs.\ extension),
la classe de risque (\texttt{risk\_class}), les scores BFPME (\texttt{score\_bfpme\_good},
\texttt{score\_bfpme\_average}) ainsi que les 62~sous-facteurs \texttt{SF\_*}. On en dérive le
score global et l'écart de score :
\[
\text{score\_bfpme\_global} =
\begin{cases}
\text{score\_bfpme\_average} & \text{si } \text{code\_modele} = 1,\\
\text{score\_bfpme\_good}    & \text{sinon,}
\end{cases}
\]
\[
\text{score\_bfpme\_gap} =
\text{score\_bfpme\_good} - \text{score\_bfpme\_average}.
\]

\subsubsection{Étape 2 --- Driver latent de risque et de qualité}

Le c\oe ur de la cohérence du jeu de données repose sur un facteur de risque latent
caché, non observable par les modèles, qui synthétise la « vraie » qualité de chaque dossier.
Il combine le score BFPME normalisé, un indice de risque dérivé de la classe de risque et le
type de projet, augmentés d'un bruit gaussien~:
\[
\begin{aligned}
\text{latent\_risk} &= \mathrm{clip}\big(r,\; 0{,}01,\; 0{,}99\big), \\[2pt]
r &= 0{,}55\Big(1-\tfrac{\text{score}}{200}\Big)
     + 0{,}35\,\text{bfpme\_risk\_index} \\
  &\quad + 0{,}10\,\mathbf{1}_{\{\text{code\_modele}=1\}}
     + \varepsilon, \\[2pt]
\varepsilon &\sim \mathcal{N}(0,\,0{,}05^2),
\end{aligned}
\]
et $\text{latent\_quality} = 1 - \text{latent\_risk}$. Sémantiquement : un meilleur score BFPME
abaisse le risque latent, une classe de risque dégradée l'augmente, et les projets de
\textit{création} sont en moyenne légèrement plus risqués que les projets d'\textit{extension}.

\subsubsection{Étape 3 --- Génération conditionnelle des variables}

Presque toutes les variables (plus de 90~au total) sont ensuite générées
conditionnellement à ce profil latent, de sorte que les entreprises solides présentent
en moyenne une meilleure liquidité, une gouvernance plus stable, des garanties plus couvrantes
et un meilleur comportement de remboursement, et inversement pour les entreprises fragiles.
Trois mécanismes sont employés~:
\begin{itemize}
  \item variables continues : moyenne pilotée par \texttt{latent\_quality} ou
        \texttt{latent\_risk} plus un bruit gaussien additif, puis écrêtage à un domaine
        métier plausible (ex.\ effectif $\in[1,250]$, ratios bornés). Exemple~:
        $\text{nombre\_employes} = \mathrm{round}(4 + 120\,\text{latent\_quality} + \mathcal{N}(0,10^2))$ ;
  \item variables binaires : tirages de Bernoulli dont la probabilité dépend du profil
        latent (ex.\ \texttt{decouverts\_frequents}, \texttt{baisse\_depots}) ;
  \item variables catégorielles : tirage multinomial pondéré (\texttt{weighted\_choice})
        dont le vecteur de probabilités varie par paliers de qualité/risque (ex.\ structure
        juridique, type de garantie, scénario macro-économique).
\end{itemize}
Les variables financières sont rendues mutuellement cohérentes (le total de bilan agrège
capitaux propres, dettes et trésorerie ; les ratios DSCR, \textit{debt-to-equity},
liquidité générale, couverture par garanties sont recalculés à partir des grandeurs générées).

\subsubsection{Étape 4 --- Bruit, stochasticité et écrêtage}

La crédibilité du jeu repose sur l'injection contrôlée d'aléa à chaque variable, ce qui évite
des lignes déterministes et reproduit l'incertitude irréductible des données réelles. Trois
sources de stochasticité coexistent~: le bruit gaussien additif ($\mathcal{N}(0,\sigma^2)$)
sur les grandeurs continues, les tirages de Bernoulli sur les indicateurs binaires, et
les tirages multinomiaux pondérés sur les variables catégorielles. Un écrêtage
systématique (\texttt{clip}) maintient chaque variable dans un domaine plausible et empêche les
valeurs aberrantes (montants négatifs, ratios extrêmes). La graine fixe garantit la
reproductibilité exacte du jeu généré.

\subsubsection{Étape 5 --- Logique de défaut synthétique}

La variable cible n'est pas tirée directement du profil latent mais d'un score de défaut
latent de forme logistique, agrégeant des pénalités et des bonus issus de variables
\emph{observables} (score, comportement de paiement, levier, DSCR, macro, protections), augmenté
d'un bruit~:
\[
\begin{aligned}
z ={} & -2{,}4
        + 2{,}5\,p_{\text{score}}
        + 1{,}6\,p_{\text{dpd}}
        + 1{,}1\,p_{\text{dscr}}
        + 1{,}0\,p_{\text{macro}}
        + 1{,}1\,p_{\text{comport.}} \\
      & - 1{,}0\,b_{\text{liquidité}}
        - 0{,}9\,b_{\text{garantie}}
        + \dots
        + \mathcal{N}(0,\,0{,}45^2),
\end{aligned}
\]
\[
\text{pd\_target} = \sigma(z) = \frac{1}{1+e^{-z}},
\qquad
\text{default\_flag} \sim \mathrm{Bernoulli}(\text{pd\_target}).
\]
Le défaut est donc un tirage stochastique et non un seuillage déterministe de la PD,
ce qui introduit naturellement des cas contre-intuitifs (de bons dossiers qui font défaut, des
dossiers risqués qui survivent), à l'image des portefeuilles réels.

\subsection{Analyse exploratoire (EDA)}

\subsubsection{Statistiques descriptives}

Examen des distributions et des valeurs aberrantes des variables, ainsi que de l'équilibre de
la cible. Le jeu final compte environ 10~000~observations et plus de 90~variables.
Le service \texttt{eda\_service} expose ce résumé (taux de défaut, complétude, nombre de
secteurs, répartitions) au tableau de bord.

\subsubsection{Distribution des défauts par secteur}

Analyse du taux de défaut selon le secteur d'activité des PME, mettant en évidence
l'hétérogénéité sectorielle reproduite par le générateur.

\subsubsection{Analyse des corrélations}

Étude des corrélations entre variables pour anticiper la multicolinéarité, notamment entre les
nombreux ratios financiers dérivés des mêmes grandeurs de base.

\subsection{Description détaillée du jeu de données et des variables}

Cette section documente le dictionnaire des variables (fichier
\texttt{dataset\_\allowbreak final\_\allowbreak variable\_\allowbreak dictionary.txt})~:
type, domaine de valeurs, logique métier de génération et transformations appliquées. Les
variables se répartissent en familles : identité d'entreprise, structure financière,
caractéristiques du crédit, garanties, environnement macro-économique, comportement de paiement
et scores BFPME, complétées par la cible \texttt{default\_flag} et la PD latente
\texttt{pd\_target}.

\subsection{Résultats du sprint}

Le sprint livre un jeu de données synthétique exploré et documenté, métier-cohérent et
reproductible, prêt à être prétraité.

\section{Sprint~3 : Prétraitement, ingénierie des variables et réalisme}

\subsection{Backlog du sprint}

Ce sprint prépare les données pour la modélisation et renforce le réalisme du jeu de données
afin d'éviter des performances artificiellement parfaites.

\subsection{Nettoyage et traitement des données}

\subsubsection{Traitement des valeurs manquantes}

Imputation selon la nature et le taux de valeurs manquantes : médiane pour les variables
numériques, modalité la plus fréquente pour les variables catégorielles. Ce choix est cohérent
avec le préprocesseur sérialisé décrit au chapitre~3.

\subsubsection{Encodage des variables catégorielles}

Transformation des variables qualitatives en représentations numériques : \texttt{OrdinalEncoder}
avec ordre métier pour les variables ordinales (ex.\ durée de crédit), \texttt{OneHotEncoder}
pour les variables nominales (secteur, structure juridique, type de garantie, scénario macro).

\subsubsection{Winsorisation, normalisation et standardisation}

Écrêtage des valeurs extrêmes par winsorisation (centiles 1 et 99) sur les grandeurs
monétaires, suivi d'une standardisation ($z$-score) des variables continues pour
homogénéiser leur influence sur les modèles sensibles à l'échelle.

\subsection{Ingénierie des variables}

\subsubsection{Indicateurs de ratio d'endettement et de solvabilité PME}

Construction de ratios financiers discriminants (\texttt{debt\_to\_equity\_ratio}, \texttt{dscr},
\texttt{ratio\_liquidite\_generale}, \texttt{quick\_ratio}, \texttt{loan\_to\_revenue},
\texttt{cash\_to\_debt}) reflétant la solvabilité et la pression d'endettement.

\subsubsection{Indicateurs sectoriels et comportementaux}

Variables synthétisant le profil sectoriel (\texttt{niveau\_risque\_sectoriel\_projet}) et le
comportement de remboursement (\texttt{payment\_stress\_index}, \texttt{impayes\_to\_loan},
\texttt{frequence\_retards}).

\subsubsection{Variables macro-économiques et prospectives}

Indicateurs intégrant une dimension prospective liée au cycle économique : variables
macro-économiques par scénario (\texttt{scenario\_macro}, \texttt{croissance\_pib},
\texttt{taux\_chomage}, \texttt{choc\_economique}) reliant le risque de défaut au contexte
conjoncturel, utiles à la construction des scénarios de stress.

\subsection{Renforcement du réalisme du jeu de données}

Une première version du jeu présentait des performances quasi parfaites (AUC proche de
0{,}99), peu crédibles pour un problème de risque de crédit réel. Une version « réalisme~v2 » a
donc été élaborée pour rapprocher l'évaluation des conditions opérationnelles, selon quatre
leviers~:
\begin{itemize}
  \item ajout de stochasticité : renforcement de la composante aléatoire dans la
        génération de la cible, augmentant l'incertitude irréductible ;
  \item réduction de la dominance des prédicteurs forts : diminution du poids des
        variables dominantes (\texttt{score\_bfpme\_global}, DPD) et répartition du signal de
        risque sur davantage de variables (macro, levier, DSCR, stress, protections) ;
  \item injection de contradictions réalistes : introduction de cas contre-intuitifs
        (bons dossiers en défaut, dossiers risqués sains) forçant les modèles à gérer le
        recouvrement entre classes ;
  \item découpage temporel : passage d'un découpage aléatoire à un découpage temporel
        (entraînement sur dossiers anciens, test sur dossiers récents) introduisant une dérive
        de distribution naturelle (\textit{time drift}).
\end{itemize}
L'effet conjugué de ces leviers ramène l'AUC à un intervalle 0{,}80--0{,}82 (Option~A),
plus plausible et défendable pour un système de risque de crédit, sans dégrader la cohérence
métier du jeu.

\subsection{Évaluation de la qualité et des biais}

\subsubsection{Analyse des biais et équité}

Détection des biais potentiels liés aux variables sensibles (localisation, structure juridique),
afin de garantir l'équité des décisions et l'absence de discrimination.

\subsubsection{Stabilité des distributions (PSI)}

Mesure de la stabilité des variables entre les sous-ensembles temporels via l'indice
\textit{Population Stability Index} (PSI), afin de quantifier la dérive introduite par le
découpage temporel.

\subsection{Résultats du sprint}

À l'issue de ce sprint, le jeu de données est entièrement nettoyé, enrichi et validé. Il
constitue une base robuste, documentée, réaliste et conforme aux exigences de qualité, prête à
alimenter la phase de modélisation présentée dans la release suivante.

\section*{Conclusion}

Cette release a permis de comprendre le métier du crédit PME et de constituer un socle de
données synthétique robuste : un générateur piloté par un facteur de risque latent, une
génération conditionnelle bruitée de plus de 90~variables métier-cohérentes, une cible obtenue
par tirage stochastique d'une PD logistique, et un renforcement explicite du réalisme
(stochasticité, réduction de dominance, contradictions, découpage temporel). La release suivante
exploite ces données pour construire et expliquer les modèles prédictifs.

% ════════════════════════════════════════════════════════════
%  CHAPITRE 5 – RELEASE II
% ════════════════════════════════════════════════════════════
\chapter{Release~II : Modélisation prédictive du risque de crédit et Explainable AI}

\section*{Introduction}

Cette release transforme le jeu de données synthétique de la release précédente en un
moteur de prédiction opérationnel. Elle poursuit deux objectifs indissociables~: d'une
part estimer la probabilité de défaut (PD) d'une PME à l'aide de modèles d'apprentissage
automatique, d'autre part rendre chaque prédiction \emph{explicable}, condition
indispensable à son adoption par les analystes crédit et à sa conformité réglementaire.
La démarche est organisée en trois sprints~: la préparation des données et du protocole
d'évaluation (Sprint~4), l'entraînement et la comparaison des modèles (Sprint~5), puis la
couche d'explicabilité et de traçabilité (Sprint~6). Le modèle finalement déployé est un
\textbf{XGBoost} entraîné sur le jeu de variables hybride de l'option~C, sérialisé avec son
préprocesseur dans un artefact unique consommé par l'API.

% ────────────────────────────────────────────────────────────
\section{Sprint~4 : Préparation des données et protocole d'évaluation}

\subsection{Backlog du sprint}

Ce sprint outille la modélisation~: ingénierie de variables discriminantes, chaîne de
prétraitement reproductible, et définition d'un protocole d'évaluation commun à tous les
modèles comparés. L'objectif est qu'un même dossier PME soit traité de façon strictement
identique à l'entraînement et à l'inférence en production.

\subsection{Ingénierie de variables}

Avant tout apprentissage, un ensemble de ratios métier est dérivé des grandeurs brutes du
dossier. Ces variables condensent des signaux de risque que les modèles exploiteraient
plus difficilement à partir des colonnes isolées~: endettement total, intensité du prêt
rapportée à l'activité, couverture de la dette par la trésorerie, pression de paiement et
couverture par les garanties. Le tableau~\ref{tab:features_engineered} en récapitule les
principales.

\begin{table}[H]
\centering
\caption{Principales variables dérivées à l'ingénierie de variables}
\label{tab:features_engineered}
\renewcommand{\arraystretch}{1.3}
\begin{tabularx}{\textwidth}{@{}l Y@{}}
\toprule
\textbf{Variable dérivée} & \textbf{Définition} \\
\midrule
\texttt{debt\_total}            & Dettes court terme + dettes long terme \\
\texttt{loan\_to\_revenue}      & Montant du prêt / chiffre d'affaires annuel \\
\texttt{cash\_to\_debt}         & Trésorerie disponible / endettement total \\
\texttt{impayes\_to\_loan}      & Montant des impayés / montant du prêt \\
\texttt{payment\_stress\_index} & Indice composite de retards et d'impayés (borné) \\
\texttt{fdr\_bfr\_pressure}     & Besoin en fonds de roulement / fonds de roulement \\
\texttt{garantie\_to\_loan}     & Valeur des garanties / montant du prêt \\
\bottomrule
\end{tabularx}
\end{table}

\noindent
Cette ingénierie est appliquée \emph{avant} le préprocesseur, de manière identique dans le
script d'entraînement et dans le service d'inférence~; elle n'est donc pas figée dans
l'artefact, ce qui garantit la cohérence entre les deux contextes.

\subsection{Chaîne de prétraitement}

Les variables sont ensuite traitées par un \texttt{ColumnTransformer} qui applique un
traitement distinct selon la nature de chaque colonne. Les variables numériques sujettes
aux valeurs extrêmes passent par une imputation médiane, un écrêtage (\textit{winsorisation}
aux 1\up{er} et 99\up{e} percentiles) puis une standardisation~; les variables numériques
simples sont imputées et standardisées~; les variables ordinales sont encodées selon un ordre
métier explicite~; les variables nominales sont encodées en \textit{one-hot}. La
figure~\ref{fig:pipeline_prepro} (release précédente) en donne la vue d'ensemble. Sur le jeu
final, ce prétraitement porte les \textbf{52~variables} d'entrée à \textbf{80~variables}
transformées.

\subsection{Protocole d'évaluation}

L'ensemble des modèles est évalué selon un protocole unique~: partition
\textit{train}/\textit{test} stratifiée à 80/20 sur la cible \texttt{default\_flag}, graine
aléatoire fixe (\texttt{random\_state~=~42}) pour la reproductibilité, et batterie de
métriques communes calculées sur l'échantillon de test~: aire sous la courbe ROC (AUC),
précision, rappel, F1-score, exactitude et matrice de confusion. Le jeu final comprend
\textbf{8~000} dossiers d'entraînement et \textbf{2~000} dossiers de test.

% ────────────────────────────────────────────────────────────
\section{Sprint~5 : Entraînement et comparaison des modèles}

\subsection{Backlog du sprint}

Entraîner et comparer les modèles selon deux axes~: le \emph{jeu de variables} retenu
(trois options de sélection) et la \emph{famille d'algorithmes} (forêt aléatoire contre
gradient boosting), afin d'arbitrer entre pouvoir discriminant, robustesse et
explicabilité métier.

\subsection{Trois options de sélection de variables}

Trois notebooks de classification explorent trois périmètres de variables, correspondant à
trois compromis entre performance pure et explicabilité réglementaire (tableau~\ref{tab:options}).

\begin{table}[H]
\centering
\caption{Les trois options de sélection de variables}
\label{tab:options}
\renewcommand{\arraystretch}{1.3}
\begin{tabularx}{\textwidth}{@{}l Y l@{}}
\toprule
\textbf{Option} & \textbf{Périmètre de variables} & \textbf{Statut} \\
\midrule
A --- \textit{All features}    & Toutes les variables utiles, scores et classe BFPME complets & Benchmark \\
B --- \textit{Without BFPME}   & Retrait de toute la famille de scores BFPME (approche \textit{data-driven}) & Benchmark \\
C --- \textit{Hybrid minimal}  & Un seul signal BFPME conservé (\texttt{score\_bfpme\_global}), 62~sous-facteurs \texttt{SF\_*} retirés & \textbf{Déployé} \\
\bottomrule
\end{tabularx}
\end{table}

\noindent
L'option~A fixe une borne supérieure de performance en tirant parti de tout le signal
disponible. L'option~B mesure la valeur prédictive propre des nouvelles variables, hors
scorecard historique. L'option~C, retenue pour le déploiement, réalise un compromis~: elle
conserve un unique signal BFPME agrégé pour préserver l'explicabilité métier et la continuité
réglementaire, tout en éliminant les 62~sous-facteurs \texttt{SF\_*} qui ajoutaient du bruit
sans gain discriminant notable.

\subsection{Familles d'algorithmes comparées}

Pour chaque option, deux familles de modèles ensemblistes à base d'arbres sont entraînées et
mises en concurrence~:
\begin{itemize}
  \item \textbf{Random Forest} --- agrégation d'arbres décorrélés par \textit{bagging},
        réputée robuste et peu sensible au surapprentissage~;
  \item \textbf{XGBoost} --- \textit{gradient boosting} d'arbres, construction séquentielle
        où chaque arbre corrige les erreurs des précédents, offrant en général le meilleur
        pouvoir discriminant sur données tabulaires.
\end{itemize}
Le tableau~\ref{tab:hyperparams} détaille la configuration retenue pour chaque modèle. Les
hyperparamètres ont été fixés à des valeurs éprouvées pour ce type de problème plutôt que
recherchés de façon exhaustive, le pouvoir discriminant étant déjà élevé sur ce jeu synthétique.

\begin{table}[H]
\centering
\caption{Configuration des modèles comparés}
\label{tab:hyperparams}
\renewcommand{\arraystretch}{1.3}
\begin{tabularx}{\textwidth}{@{}l Y@{}}
\toprule
\textbf{Modèle} & \textbf{Hyperparamètres principaux} \\
\midrule
Random Forest & \texttt{n\_estimators=400}, \texttt{min\_samples\_split=8}, \texttt{min\_samples\_leaf=4} \\
XGBoost       & \texttt{n\_estimators=500}, \texttt{max\_depth=5}, \texttt{learning\_rate=0{.}05}, \texttt{subsample=0{.}85}, \texttt{colsample\_bytree=0{.}85}, \texttt{objective=binary:logistic} \\
\bottomrule
\end{tabularx}
\end{table}

\subsection{Résultats comparatifs}

Les courbes ROC et les métriques comparées des options A, B et C
(figure~\ref{fig:roc_comparaison}) montrent des performances élevées et proches entre les
trois périmètres, avec des écarts faibles entre Random Forest et XGBoost. Ce comportement est
attendu sur un jeu synthétique dont la cible est dérivée d'un score latent réutilisant des
signaux présents en entrée~: les performances valident la cohérence technique du pipeline
mais ne préjugent pas d'un niveau réaliste sur données réelles (voir la discussion de la
release~I sur l'injection de réalisme).

\begin{figure}[H]
  \centering
  \begin{tikzpicture}[x=5.8cm, y=5.8cm]
    % grille
    \foreach \g in {0.2,0.4,0.6,0.8}{
      \draw[very thin, gray!25] (\g,0) -- (\g,1);
      \draw[very thin, gray!25] (0,\g) -- (1,\g);
    }
    % axes
    \draw[->, thick] (0,0) -- (1.06,0)
      node[right, font=\footnotesize] {Taux de faux positifs};
    \draw[->, thick] (0,0) -- (0,1.06)
      node[above, font=\footnotesize] {Taux de vrais positifs};
    \foreach \t in {0,0.2,0.4,0.6,0.8,1}{
      \draw (\t,0) -- (\t,-0.018) node[below, font=\tiny] {\t};
      \draw (0,\t) -- (-0.018,\t) node[left,  font=\tiny] {\t};
    }
    % diagonale (classifieur aléatoire)
    \draw[dashed, gray] (0,0) -- (1,1);
    % courbes ROC
    \draw[very thick, EspritRed]
      plot[smooth] coordinates {(0,0) (0.04,0.45) (0.10,0.66) (0.20,0.80)
                                (0.35,0.905) (0.55,0.96) (0.80,0.99) (1,1)};
    \draw[very thick, DarkGray]
      plot[smooth] coordinates {(0,0) (0.05,0.40) (0.12,0.61) (0.22,0.77)
                                (0.40,0.89) (0.60,0.95) (0.82,0.985) (1,1)};
    \draw[very thick, gray!55]
      plot[smooth] coordinates {(0,0) (0.06,0.35) (0.15,0.56) (0.27,0.73)
                                (0.47,0.865) (0.67,0.935) (0.86,0.98) (1,1)};
    % légende
    \begin{scope}[shift={(0.34,0.05)}]
      \fill[white, opacity=0.82] (-0.02,-0.01) rectangle (0.62,0.30);
      \draw[very thick, EspritRed] (0,0.24)--(0.08,0.24)
        node[right, font=\tiny, text=black] {Option C (XGBoost) --- AUC 0,89};
      \draw[very thick, DarkGray]  (0,0.14)--(0.08,0.14)
        node[right, font=\tiny, text=black] {Option A --- AUC 0,88};
      \draw[very thick, gray!55]   (0,0.04)--(0.08,0.04)
        node[right, font=\tiny, text=black] {Option B --- AUC 0,86};
    \end{scope}
  \end{tikzpicture}
  \caption{Courbes ROC comparées des options A, B et C (Random Forest vs.\ XGBoost)}
  \label{fig:roc_comparaison}
\end{figure}

\subsection{Sélection du modèle déployé}

Le modèle retenu est le \textbf{XGBoost de l'option~C}, qui offre le meilleur compromis entre
pouvoir discriminant, parcimonie et explicabilité métier. Le tableau~\ref{tab:metrics_deploy}
reporte ses performances sur l'échantillon de test.

\begin{table}[H]
\centering
\caption{Performances du modèle déployé (XGBoost, option~C) sur l'échantillon de test}
\label{tab:metrics_deploy}
\renewcommand{\arraystretch}{1.3}
\begin{tabularx}{\textwidth}{@{}l c Y@{}}
\toprule
\textbf{Métrique} & \textbf{Valeur} & \textbf{Interprétation} \\
\midrule
AUC-ROC   & 0{,}889 & Pouvoir discriminant (Gini $= 2\,\text{AUC}-1 \approx 0{,}78$) \\
Exactitude& 0{,}855 & Proportion de prédictions correctes \\
Précision & 0{,}868 & Fiabilité des défauts prédits \\
Rappel    & 0{,}875 & Couverture des défauts réels \\
F1-score  & 0{,}872 & Moyenne harmonique précision/rappel \\
\bottomrule
\end{tabularx}
\end{table}

\noindent
Le préprocesseur ajusté et le modèle entraîné sont sérialisés ensemble (\texttt{joblib}) dans
un artefact unique, accompagné de la liste des colonnes autorisées~; cet artefact est chargé
et mis en cache au démarrage de l'API, garantissant une inférence strictement alignée sur
l'entraînement.

% ────────────────────────────────────────────────────────────
\section{Sprint~6 : Explainable AI (XAI) et traçabilité}

\subsection{Backlog du sprint}

Doter le modèle déployé d'une couche d'explicabilité complète, produisant pour chaque
prédiction des explications globales (à l'échelle du portefeuille) et locales (pour un
dossier donné), dans un format auditable par les analystes et les régulateurs.

\subsection{Explicabilité globale}

L'explicabilité globale repose sur les valeurs de Shapley (SHAP), calculées par un
\texttt{TreeExplainer} directement sur la matrice de variables transformées. Le
\textit{summary plot} (figure~\ref{fig:shap_summary}) classe les variables par contribution
moyenne absolue et révèle le sens de leur effet sur le risque à l'échelle du portefeuille. En
complément, une \textit{importance par permutation} est calculée pour le Random Forest sur les
colonnes d'origine, offrant une lecture métier directe, exprimée dans les variables non
encodées.

\begin{figure}[H]
  \centering
  \begin{tikzpicture}[x=1cm, y=1cm]
    \def\cx{5.4}   % abscisse de la ligne SHAP = 0
    \def\sc{9.5}   % échelle : cm par unité SHAP
    % lignes de base + libellés des variables (ordre = importance décroissante)
    \foreach \i in {1,...,6}{\draw[very thin, gray!25] (1.6,{-0.72*\i}) -- (9.6,{-0.72*\i});}
    \node[left, font=\scriptsize\ttfamily] at (1.5,-0.72) {score\_bfpme\_global};
    \node[left, font=\scriptsize\ttfamily] at (1.5,-1.44) {payment\_stress\_index};
    \node[left, font=\scriptsize\ttfamily] at (1.5,-2.16) {cash\_to\_debt};
    \node[left, font=\scriptsize\ttfamily] at (1.5,-2.88) {loan\_to\_revenue};
    \node[left, font=\scriptsize\ttfamily] at (1.5,-3.60) {garantie\_to\_loan};
    \node[left, font=\scriptsize\ttfamily] at (1.5,-4.32) {retards\_paiement};
    % ligne SHAP = 0
    \draw[dashed, gray] (\cx,-0.35) -- (\cx,-4.75);
    \node[font=\tiny, above] at (\cx,-0.30) {SHAP $= 0$};
    % nuages de points (position = impact ; gauche/droite = sens)
    \foreach \v/\j in {-0.16/0.05, -0.05/-0.07, 0.06/0.09, 0.13/-0.03, 0.19/0.07,
                       0.25/-0.06, 0.32/0.04, 0.10/0.12, 0.21/-0.11}
      \fill[DarkGray, opacity=0.55] ({\cx+\v*\sc}, {-0.72*1+\j}) circle (1.7pt);
    \foreach \v/\j in {0.02/0.06, 0.07/-0.05, 0.12/0.10, 0.17/0.0, 0.22/-0.08,
                       0.27/0.07, 0.15/-0.11, 0.09/0.12}
      \fill[DarkGray, opacity=0.55] ({\cx+\v*\sc}, {-0.72*2+\j}) circle (1.7pt);
    \foreach \v/\j in {-0.24/0.05, -0.18/-0.06, -0.12/0.08, -0.07/-0.02, -0.02/0.09,
                       0.03/-0.07, -0.15/0.11, -0.09/-0.10}
      \fill[DarkGray, opacity=0.55] ({\cx+\v*\sc}, {-0.72*3+\j}) circle (1.7pt);
    \foreach \v/\j in {0.0/0.06, 0.04/-0.05, 0.08/0.10, 0.12/0.0, 0.16/-0.08,
                       0.20/0.07, 0.06/-0.11, 0.14/0.12}
      \fill[DarkGray, opacity=0.55] ({\cx+\v*\sc}, {-0.72*4+\j}) circle (1.7pt);
    \foreach \v/\j in {-0.19/0.05, -0.14/-0.06, -0.09/0.08, -0.05/-0.02, -0.01/0.09,
                       0.02/-0.07, -0.11/0.11, -0.07/-0.10}
      \fill[DarkGray, opacity=0.55] ({\cx+\v*\sc}, {-0.72*5+\j}) circle (1.7pt);
    \foreach \v/\j in {0.0/0.06, 0.03/-0.05, 0.06/0.10, 0.09/0.0, 0.12/-0.08,
                       0.15/0.07, 0.05/-0.11, 0.10/0.12}
      \fill[DarkGray, opacity=0.55] ({\cx+\v*\sc}, {-0.72*6+\j}) circle (1.7pt);
    % annotations de sens
    \node[font=\scriptsize\itshape, anchor=east] at (\cx-0.3,-5.15)
      {$\leftarrow$ diminue le risque};
    \node[font=\scriptsize\itshape, anchor=west] at (\cx+0.3,-5.15)
      {augmente le risque $\rightarrow$};
  \end{tikzpicture}
  \caption{SHAP \textit{summary plot}~: importance et direction des variables (portefeuille)}
  \label{fig:shap_summary}
\end{figure}

\subsection{Explicabilité locale}

Pour un dossier PME donné, le service d'explication restitue la valeur de base
(\textit{expected value}), la prédiction, et les contributions SHAP ordonnées des variables
les plus influentes, chacune étiquetée selon sa direction (\texttt{increase\_risk} ou
\texttt{decrease\_risk}). Le \textit{waterfall plot} (figure~\ref{fig:shap_waterfall})
décompose ainsi le passage de la valeur de base à la PD prédite, variable par variable.

\begin{figure}[H]
  \centering
  \begin{tikzpicture}[x=1cm, y=1cm]
    % Échelle : x(v) = 1 + (v - 0.25) * 20  (PD de 0,25 à 0,70)
    % Contributions cumulées :
    %  base 0,30 -> +0,13 -> +0,10 -> +0,07 -> -0,05 -> -0,04 -> +0,11 = 0,62
    % axe des PD (bas)
    \draw[->, thick] (1,-5.1) -- (10.2,-5.1) node[right, font=\scriptsize] {PD};
    \foreach \v in {0.30,0.40,0.50,0.60}{
      \pgfmathsetmacro\xx{1+(\v-0.25)*20}
      \draw (\xx,-5.1) -- (\xx,-5.22) node[below, font=\tiny] {\v};
    }
    % lignes de référence : valeur de base et prédiction
    \draw[dashed, gray] (2.0,-0.3) -- (2.0,-5.1);
    \node[font=\tiny, gray, above] at (2.0,-0.3) {$E[f(x)]=0{,}30$};
    \draw[dashed, EspritRed] (8.4,-0.3) -- (8.4,-5.1);
    \node[font=\tiny, EspritRed, above] at (8.4,-0.3) {$f(x)=0{,}62$};
    % barres flottantes (increase = rouge, decrease = gris clair)
    % row i : y = -0.7*i, hauteur 0.44
    \newcommand{\wbar}[6]{% xa xb y fill label value
      \fill[#4] (#1,{#3-0.22}) rectangle (#2,{#3+0.22});
      \node[left, font=\scriptsize\ttfamily] at (1.7,#3) {#5};
      \node[right, font=\tiny] at ({max(#1,#2)+0.1},#3) {#6};}
    \wbar{2.0}{4.6}{-0.7}{EspritRed!75}{payment\_stress\_index}{$+0{,}13$}
    \wbar{4.6}{6.6}{-1.4}{EspritRed!75}{score\_bfpme\_global}{$+0{,}10$}
    \wbar{6.6}{8.0}{-2.1}{EspritRed!75}{loan\_to\_revenue}{$+0{,}07$}
    \wbar{8.0}{7.0}{-2.8}{gray!45}{cash\_to\_debt}{$-0{,}05$}
    \wbar{7.0}{6.2}{-3.5}{gray!45}{garantie\_to\_loan}{$-0{,}04$}
    \wbar{6.2}{8.4}{-4.2}{EspritRed!75}{retards\_paiement}{$+0{,}11$}
    % connecteurs (fil cumulatif)
    \foreach \x/\ya/\yb in {4.6/-0.7/-1.4, 6.6/-1.4/-2.1, 8.0/-2.1/-2.8,
                            7.0/-2.8/-3.5, 6.2/-3.5/-4.2}{
      \draw[dotted, gray] (\x,{\ya-0.22}) -- (\x,{\yb+0.22});
    }
  \end{tikzpicture}
  \caption{SHAP \textit{waterfall plot}~: décomposition d'une décision individuelle}
  \label{fig:shap_waterfall}
\end{figure}

\subsection{Format et auditabilité des explications}

Chaque explication est structurée (valeur de base, contributions ordonnées, direction, version
du modèle) de manière à être à la fois lisible par un analyste crédit — qui y lit les facteurs
concrets augmentant ou diminuant le risque du dossier — et traçable pour un régulateur — qui
dispose d'une décomposition documentée, reproductible et rattachée à une version de modèle
identifiée. Le listing~\ref{lst:shap} illustre le c\oe ur du service d'explication.

\begin{codeblock}
\begin{lstlisting}[caption={Calcul des contributions SHAP locales (extrait)}, label={lst:shap},
                   language=Python]
explainer   = TreeExplainer(model)               # sur le modele XGBoost
shap_values = explainer.shap_values(transformed) # matrice transformee (80 var.)
top_idx     = np.argsort(np.abs(shap_row))[::-1][:top_k]  # variables dominantes
# chaque contribution est etiquetee selon sa direction
direction = "increase_risk" if shap_row[idx] >= 0 else "decrease_risk"
\end{lstlisting}
\end{codeblock}

\section*{Conclusion}

Cette release livre un modèle prédictif de la PME performant et pleinement explicable~: un
XGBoost entraîné sur le jeu de variables hybride de l'option~C, comparé à une forêt aléatoire
et à deux périmètres de variables concurrents, puis doté d'une couche SHAP produisant des
explications globales et locales auditables. Le modèle et son préprocesseur sont sérialisés
dans un artefact unique, prêt à être orchestré par l'agent IA de la release suivante.

% ════════════════════════════════════════════════════════════
%  CHAPITRE 6 – RELEASE III  (Agent IA + LLM)
% ════════════════════════════════════════════════════════════
\chapter{Release~III : Agent IA et intégration LLM}

\section*{Introduction}

Cette release industrialise le modèle prédictif de la release~II au sein d'un \emph{agent IA
orchestrateur}, puis lui adjoint une couche d'interaction en langage naturel. L'agent
matérialise la logique de décision de la plateforme~: à partir d'un dossier PME et d'une
requête d'analyste, il décide quels traitements invoquer (validation, prédiction, explication,
simulation de scénario), compose leurs résultats en une réponse unique, et — si elle est
activée — délègue à un grand modèle de langage (LLM) la seule production d'une synthèse
narrative. Le principe directeur, gravé dans l'implémentation, est une \textbf{séparation
stricte des rôles}~: les valeurs numériques sont exclusivement calculées par le modèle et ses
services, jamais par le LLM. Concrètement, l'agent est exposé par le point d'entrée
\texttt{POST /chat/analyze}.

\section{Sprint~7 : Architecture de l'agent IA orchestrateur}

\subsection{Backlog du sprint}

Concevoir et implémenter l'agent orchestrateur~: enchaînement conditionnel des outils,
composition des résultats, robustesse aux pannes de la couche LLM, et classes de risque pour
l'escalade.

\subsection{Pipeline d'orchestration}

L'agent est implémenté comme un routeur asynchrone (\texttt{routers/chat.py}) qui, à réception
d'un dossier, invoque \emph{sélectivement} les services métier selon des drapeaux de requête,
et accumule leurs sorties dans un dictionnaire \texttt{tool\_results}. Chaque service est
découplé de la couche HTTP et réutilisable~; \texttt{model\_service.predict} est ainsi appelé
à la fois par la prédiction directe, par le scénario et par l'orchestration. Le
listing~\ref{lst:orchestration} montre le c\oe ur de cet enchaînement.

\begin{codeblock}
\begin{lstlisting}[caption={Enchaînement conditionnel des outils par l'agent}, label={lst:orchestration},
                   language=Python]
if payload.run_validation:
    tool_results["validation"]  = validate_features(features)
if payload.run_prediction:
    tool_results["prediction"]  = predict(features, threshold)
if payload.run_explanation:
    tool_results["explanation"] = explain_prediction(features, top_k_shap)
if payload.scenario_changes:
    tool_results["scenario"]    = run_scenario(features, changes, threshold)
# la synthese narrative n'intervient qu'ensuite, et seulement si le LLM est actif
\end{lstlisting}
\end{codeblock}

\noindent
Le tableau~\ref{tab:tools_agent} récapitule les outils que l'agent peut mobiliser et le
service qui les réalise.

\begin{table}[H]
\centering
\caption{Outils orchestrés par l'agent IA}
\label{tab:tools_agent}
\renewcommand{\arraystretch}{1.3}
\begin{tabularx}{\textwidth}{@{}l l Y@{}}
\toprule
\textbf{Outil} & \textbf{Service} & \textbf{Rôle} \\
\midrule
Validation  & \texttt{validation\_service} & Vérifie les variables attendues, signale manquantes / non utilisées \\
Prédiction  & \texttt{model\_service}      & Estime la PD et attribue la classe de risque \\
Explication & \texttt{shap\_service}       & Calcule les contributions SHAP locales \\
Scénario    & \texttt{scenario\_service}   & Rejoue le modèle sous une hypothèse \textit{what-if} \\
Synthèse    & \texttt{llm\_service}        & Rédige un récit à partir des seules sorties calculées \\
\bottomrule
\end{tabularx}
\end{table}

\subsection{Classes de risque et seuils}

L'agent distingue deux notions. D'une part un \emph{seuil de décision} configurable
(\texttt{DEFAULT\_THRESHOLD}, 0{,}5 par défaut)~: la classe prédite vaut « défaut » dès que la
PD atteint ce seuil. D'autre part une \emph{classe de risque qualitative}, dérivée
directement de la PD, qui sert de support à l'escalade~:
\[
\text{classe} =
\begin{cases}
\text{\textit{high}}   & \text{si } \text{PD} \geq 0{,}70,\\
\text{\textit{medium}} & \text{si } 0{,}40 \leq \text{PD} < 0{,}70,\\
\text{\textit{low}}    & \text{sinon.}
\end{cases}
\]
Un dossier classé \textit{high} ou proche du seuil constitue naturellement un cas d'escalade
vers le comité de crédit, accompagné de ses explications SHAP — la plateforme autonomise
l'analyste sans jamais automatiser la décision finale.

\subsection{Robustesse et résilience}

L'agent est conçu pour que la couche LLM ne puisse jamais compromettre les résultats calculés.
Si le LLM est désactivé, l'agent renvoie les sorties brutes des outils. S'il est activé mais
qu'une erreur survient (serveur injoignable, délai dépassé), l'exception est capturée~: les
\texttt{tool\_results} sont préservés, le drapeau \texttt{used\_llm} passe à \texttt{False}, et
un message diagnostique explicite est retourné. La valeur métier — PD, explications, scénario —
est donc toujours délivrée, la synthèse narrative n'étant qu'une couche de confort.

\subsection{Simulation de scénarios \textit{what-if}}

L'agent peut rejouer le modèle sous une hypothèse de l'analyste. Le \texttt{scenario\_service}
applique les changements demandés au dossier de base, relance une prédiction complète, et
retourne la PD de référence, la PD après modification, leur écart (\texttt{delta\_pd}, positif
si le risque augmente) et les classes de risque associées. Il s'agit d'un \emph{vrai}
ré-entraînement de la prédiction, non d'une estimation~: la modification est réellement propagée
dans le pipeline.

\section{Sprint~8 : Intégration LLM et interaction en langage naturel}

\subsection{Backlog du sprint}

Intégrer un LLM comme couche de communication~: traduire les sorties numériques en une synthèse
lisible et permettre à l'analyste d'interroger le système en langage naturel, sans jamais
laisser le LLM produire de valeur numérique.

\subsection{Choix du modèle et souveraineté des données}

Le LLM est servi par une infrastructure \textbf{Ollama} exposant une API compatible OpenAI,
interrogée en local. Ce choix répond aux contraintes du secteur bancaire~: les données de
crédit ne quittent pas l'infrastructure de la banque (souveraineté et confidentialité), et le
coût d'inférence est maîtrisé. Le modèle par défaut est \texttt{qwen2.5:7b-instruct}, appelé
avec une température basse (0{,}1) pour minimiser la variabilité, et un plafond de jetons
configurable. Toute la configuration (activation, URL, modèle, clé, plafond de jetons) est
pilotée par variables d'environnement.

\subsection{Séparation stricte des rôles : le LLM n'estime jamais le risque}

C'est la garantie centrale de l'intégration. Le LLM ne reçoit que le dictionnaire
\texttt{tool\_results} déjà calculé et une consigne système lui interdisant explicitement de
fabriquer, arrondir ou estimer une valeur numérique. Le \textit{prompt} lui impose de
n'utiliser que les sorties fournies et de restituer les chiffres exacts — notamment, en
présence d'un bloc scénario, la PD de référence, la PD modifiée et leur écart. La réponse est
structurée en sections imposées~: \textit{Decision Summary}, \textit{Risk Drivers},
\textit{Protective Drivers}, \textit{Scenario Impact} (uniquement si un scénario existe) et
\textit{Caveats}. Le listing~\ref{lst:prompt} en montre l'ossature.

\begin{codeblock}
\begin{lstlisting}[caption={Consigne imposée au LLM (extrait du prompt)}, label={lst:prompt}]
You are a credit risk analyst assistant for BFPME (SME lending).
Use ONLY the tool outputs below. Do not invent, round, or estimate
any numeric value.
If a 'scenario' block is present, report the EXACT numbers:
  - baseline_pd, scenario_pd, delta_pd, and the changed_features.
Structure: Decision Summary, Risk Drivers, Protective Drivers,
           Scenario Impact (only if present), Caveats.
\end{lstlisting}
\end{codeblock}

\subsection{Interface de questions-réponses pour les analystes}

Depuis le tableau de bord, l'analyste soumet une question en langage naturel accompagnée du
dossier PME. L'agent exécute les outils pertinents puis, si le LLM est actif, produit une
synthèse narrative qui traduit les contributions SHAP en énoncés compréhensibles (facteurs
aggravant ou atténuant le risque), justifie la décision et, le cas échéant, commente l'impact
d'un scénario \textit{what-if}. L'analyste peut ainsi explorer les déterminants du risque et des
hypothèses alternatives sans manipuler directement les sorties techniques du modèle.

\subsection{Gestion des pannes et délais}

L'appel au LLM est asynchrone (\texttt{httpx}) et non bloquant, avec des délais généreux en
lecture (jusqu'à 300~s) pour absorber le chargement à froid d'un modèle de 7~milliards de
paramètres. Les erreurs typiques sont traduites en messages actionnables~: serveur Ollama
injoignable, ou modèle en cours de chargement. Conformément au principe de robustesse, aucune
de ces situations ne fait perdre les résultats déjà calculés.

\section*{Conclusion}

Cette release livre un agent IA interactif et explicable~: un orchestrateur qui enchaîne
sélectivement validation, prédiction, explication et simulation de scénarios, puis délègue à un
LLM local — sous séparation stricte des rôles — la seule mise en récit des résultats. La
robustesse de l'agent garantit que la valeur métier est toujours délivrée, indépendamment de la
disponibilité du LLM. La dernière release traite des tests de stress du portefeuille par
simulation Monte Carlo.

% ════════════════════════════════════════════════════════════
%  CHAPITRE 7 – RELEASE IV
% ════════════════════════════════════════════════════════════
\chapter{Release~IV : Tests de stress par simulation Monte Carlo}

\section*{Introduction}

Cette release complète la plateforme par un moteur de \emph{tests de stress} fondé sur la
simulation stochastique de Monte Carlo. Là où l'agent de la release précédente évalue un
dossier sous une hypothèse ponctuelle (\textit{what-if}), ce moteur explore un \emph{grand
nombre} de scénarios perturbés simultanément, afin de reconstituer la distribution complète de
la probabilité de défaut (PD) d'un dossier soumis à des chocs et d'en dériver des indicateurs
de risque extrême. Il est exposé par le point d'entrée \texttt{POST /monte-carlo} et réutilise
exactement le modèle déployé, garantissant la cohérence avec la prédiction nominale.

\section{Sprint~9 : Moteur de tests de stress Monte Carlo}

\subsection{Backlog du sprint}

Mettre en place un moteur de simulation Monte Carlo vectorisé qui, à partir d'un dossier de
base et d'un ensemble de chocs sur les variables clés, génère des milliers de scénarios,
évalue la PD de chacun via le modèle déployé, et synthétise la distribution obtenue en
indicateurs de risque exploitables (VaR, Expected Shortfall, répartition par classe de risque).

\subsection{Fondements : la méthode de Monte Carlo}

La méthode de Monte Carlo estime la distribution d'une grandeur de sortie en propageant, à
travers un modèle, un grand nombre de tirages aléatoires de ses entrées~\cite{glasserman2004,
metropolis1949}. Appliquée au risque de crédit, elle consiste à perturber aléatoirement les
variables d'un dossier PME — selon des chocs représentatifs d'un environnement adverse — puis à
observer comment la PD prédite se déforme. Elle fournit ainsi non pas une valeur unique, mais
toute une distribution de PD, dont on extrait les comportements de queue. Cette démarche
s'inscrit dans les principes de stress testing du Comité de Bâle~\cite{bcbs2018}.

\subsection{Génération des scénarios perturbés}

Chaque choc porte sur une variable numérique du dossier et est défini par une amplitude
relative (\texttt{pct}). Pour $n$ simulations, le moteur construit $n+1$ lignes~: la ligne~0
est le dossier de base intact (référence), les lignes~1 à~$n$ sont des perturbations. Pour
chaque variable choquée de valeur de base $b$, un facteur multiplicatif aléatoire est tiré selon
la loi choisie~:
\[
\text{facteur} \sim
\begin{cases}
\mathcal{N}(1,\ \texttt{pct}^2) & \text{loi normale,}\\
\mathcal{U}(1-\texttt{pct},\ 1+\texttt{pct}) & \text{loi uniforme,}
\end{cases}
\qquad
\text{valeur} = b \times \text{facteur}.
\]
Les grandeurs financières positives (montants, chiffre d'affaires) sont écrêtées à zéro pour
interdire des valeurs négatives aberrantes, tandis que les taux peuvent changer de signe. Point
essentiel de cohérence~: après application des chocs sur les variables \emph{brutes},
l'ingénierie de variables est \emph{re-dérivée}, de sorte que les ratios
(\texttt{loan\_to\_revenue}, \texttt{cash\_to\_debt}, etc.) restent cohérents avec les valeurs
perturbées.

\subsection{Paramétrage de la simulation}

La requête adressée à \texttt{/monte-carlo} expose un ensemble de paramètres bornés, validés par
le schéma Pydantic \texttt{MonteCarloRequest}, qui donnent à l'analyste un contrôle fin sur le
scénario de stress (tableau~\ref{tab:mc_params}). Chaque choc est décrit par le couple
(\texttt{field}, \texttt{pct}), où \texttt{pct} est l'amplitude relative de la perturbation
(par exemple \texttt{pct~=~0{,}2} pour un choc de $\pm 20\,\%$).

\begin{table}[H]
\centering
\caption{Paramètres du moteur de simulation Monte Carlo}
\label{tab:mc_params}
\renewcommand{\arraystretch}{1.3}
\begin{tabularx}{\textwidth}{@{}l l Y@{}}
\toprule
\textbf{Paramètre} & \textbf{Valeur / plage} & \textbf{Rôle} \\
\midrule
\texttt{shocks}        & liste (\texttt{field}, \texttt{pct}$\in[0;2]$) & Variables choquées et amplitude relative \\
\texttt{n\_simulations}& 1\,000 (10 à 20\,000)          & Nombre de scénarios générés \\
\texttt{distribution}  & \texttt{normal} / \texttt{uniform} & Loi de tirage des facteurs de choc \\
\texttt{threshold}     & 0{,}5 par défaut               & Seuil de décision pour $P(\text{PD}\geq\text{seuil})$ \\
\texttt{seed}          & entier optionnel              & Graine aléatoire (reproductibilité) \\
\texttt{n\_bins}       & 10 (4 à 50)                   & Nombre de classes de l'histogramme \\
\texttt{n\_extremes}   & 3 (0 à 20)                    & Nombre de scénarios extrêmes restitués \\
\bottomrule
\end{tabularx}
\end{table}

\subsection{Évaluation vectorisée}

L'ensemble des $n+1$ scénarios est évalué en \emph{un seul passage vectorisé}~: une unique
transformation par le préprocesseur suivie d'un unique appel à \texttt{predict\_proba} sur tout
le lot. Ce choix d'implémentation rend le moteur capable de traiter des milliers de scénarios
en une fraction du temps qu'exigerait une boucle dossier par dossier. La PD de la ligne~0 sert
de référence~; les PD des lignes suivantes forment l'échantillon simulé. Le listing~\ref{lst:mc}
résume le c\oe ur de l'algorithme.

\begin{codeblock}
\begin{lstlisting}[caption={Cœur du moteur Monte Carlo (extrait)}, label={lst:mc},
                   language=Python]
# 1) tirage des facteurs de choc (loi normale ou uniforme)
if distribution == "normal":
    factors = rng.normal(1.0, shock.pct, n)
else:
    factors = rng.uniform(1.0 - shock.pct, 1.0 + shock.pct, n)
values = base[shock.field] * factors
values = np.clip(values, 0.0, None)          # grandeurs positives

# 2) re-derivation des ratios, puis evaluation VECTORISEE (un seul appel)
frame       = apply_feature_engineering(frame)
transformed = preprocessor.transform(frame[allowed_columns])
pd_all      = model.predict_proba(transformed)[:, 1]

# 3) indicateurs de queue
var_95 = np.percentile(sims, 95)
es_95  = sims[sims >= var_95].mean()          # Expected Shortfall
\end{lstlisting}
\end{codeblock}

\subsection{Validation des chocs et cohérence}

Le moteur applique plusieurs garde-fous garantissant des scénarios exploitables. Seuls les chocs
portant sur une variable \emph{numérique effectivement présente} dans le dossier sont retenus~;
les champs non numériques ou absents sont ignorés et signalés dans le champ \texttt{notes} de la
réponse, sans interrompre la simulation. Si aucun choc n'est actif, toutes les simulations
coïncident avec le dossier de base, ce qui est explicitement indiqué. La graine aléatoire
(\texttt{seed}) rend chaque simulation exactement reproductible, condition indispensable à
l'auditabilité d'un test de stress. Enfin, la ré-application de l'ingénierie de variables après
les chocs assure que les ratios dérivés restent cohérents avec les grandeurs perturbées, évitant
des scénarios financièrement contradictoires.

\subsection{Indicateurs de risque dérivés}

De la distribution des PD simulées, le moteur dérive une batterie d'indicateurs
(tableau~\ref{tab:mc_indicateurs}) synthétisant la tendance centrale, la dispersion et surtout
le risque de queue.

\begin{table}[H]
\centering
\caption{Indicateurs de risque produits par le moteur Monte Carlo}
\label{tab:mc_indicateurs}
\renewcommand{\arraystretch}{1.3}
\begin{tabularx}{\textwidth}{@{}l Y@{}}
\toprule
\textbf{Indicateur} & \textbf{Définition} \\
\midrule
PD moyenne, écart-type & Tendance centrale et dispersion de la PD sous stress \\
Percentiles (p5 \dots p95) & Forme de la distribution simulée \\
VaR 95\,\% & Percentile 95 de la PD (perte de crédit à 95\,\% de confiance) \\
Expected Shortfall 95\,\% & PD moyenne au-delà de la VaR (sévérité de la queue) \\
$P(\text{PD} \geq \text{seuil})$ & Probabilité de franchir le seuil de décision sous stress \\
Répartition \textit{low/medium/high} & Mix de classes de risque parmi les scénarios \\
Scénarios extrêmes & Dossiers les plus/moins risqués et leurs chocs appliqués \\
\bottomrule
\end{tabularx}
\end{table}

\noindent
La \textbf{VaR à 95\,\%} correspond au 95\up{e} percentile de la distribution des PD~: dans
95\,\% des scénarios, la PD reste en deçà de cette valeur. L'\textbf{Expected Shortfall} (ou
\textit{CVaR}) est la moyenne des PD au-delà de la VaR~; elle mesure la sévérité \emph{dans} la
queue, information que la VaR seule ne capture pas. Le moteur restitue également un histogramme
de la distribution et les scénarios extrêmes (les plus et les moins risqués), assortis des
valeurs choquées qui les ont produits, ce qui permet à l'analyste d'identifier les combinaisons
de chocs les plus pénalisantes.

\subsection{Analyse de la distribution et robustesse}

La figure~\ref{fig:mc_distribution} illustre une distribution de PD simulée~: le déplacement de
sa masse et l'épaisseur de sa queue droite, comparés à la PD de référence, renseignent
directement sur la robustesse du dossier face aux chocs. Une queue droite épaisse — VaR et
Expected Shortfall nettement supérieures à la PD de référence — signale un dossier fragile,
dont le risque se dégrade fortement sous stress~; à l'inverse, une distribution resserrée
autour de la référence traduit une bonne résilience.

\begin{figure}[H]
  \centering
  \begin{tikzpicture}[x=11.5cm, y=5cm]
    % ── cadre et axes ────────────────────────────────────────
    \draw[->, thick] (0,0) -- (1.04,0) node[right, font=\footnotesize] {PD};
    \draw[->, thick] (0,0) -- (0,1.06) node[above, font=\footnotesize] {Fréquence};
    \foreach \t in {0,0.2,0.4,0.6,0.8,1.0}{
      \draw (\t,0) -- (\t,-0.02) node[below, font=\tiny] {\t};
    }
    % ── histogramme des PD simulées (hauteurs normalisées) ───
    % centres de classe / hauteurs (distribution décalée vers la droite, queue épaisse)
    \foreach \c/\h in {0.05/0.06, 0.15/0.16, 0.25/0.34, 0.35/0.58, 0.45/0.82,
                       0.55/1.00, 0.65/0.80, 0.75/0.52, 0.85/0.30, 0.95/0.14}{
      \fill[gray!35, draw=gray!60] ({\c-0.05},0) rectangle ({\c+0.05},\h);
    }
    % ── PD de référence (baseline) ──────────────────────────
    \draw[thick, DarkGray, dashed] (0.30,0) -- (0.30,1.02);
    \node[above, font=\tiny, text=DarkGray, align=center] at (0.30,1.02)
      {PD de\\référence\\0,30};
    % ── VaR 95 % ────────────────────────────────────────────
    \draw[very thick, EspritRed] (0.78,0) -- (0.78,0.92);
    \node[above, font=\scriptsize\bfseries, text=EspritRed] at (0.78,0.92) {VaR 95\,\% = 0,68};
    % ── zone de queue (au-delà de la VaR) + Expected Shortfall ─
    \foreach \c/\h in {0.85/0.30, 0.95/0.14}{
      \fill[EspritRed!25, draw=EspritRed!50] ({\c-0.05},0) rectangle ({\c+0.05},\h);
    }
    \draw[very thick, EspritRed, densely dotted] (0.90,0) -- (0.90,0.55);
    \node[above right, font=\scriptsize, text=EspritRed] at (0.90,0.42)
      {\; ES 95\,\% = 0,77};
    % ── légende de la queue ─────────────────────────────────
    \node[font=\tiny, text=EspritRed, align=center] at (0.90,-0.10)
      {5\,\% des scénarios\\les plus adverses};
  \end{tikzpicture}
  \caption{Distribution des PD simulées sous chocs (VaR et Expected Shortfall à 95\,\%)}
  \label{fig:mc_distribution}
\end{figure}

\subsection{Exemple de lecture pour l'analyste}

Concrètement, l'analyste applique par exemple un choc de $-20\,\%$ sur le chiffre d'affaires et
de $+15\,\%$ sur les impayés d'un dossier, puis lit la réponse du moteur. Une PD de référence de
$0{,}30$ qui se déplace vers une PD moyenne de $0{,}42$, avec une VaR à 95\,\% de $0{,}68$ et une
Expected Shortfall de $0{,}77$, indique un dossier dont le risque, initialement modéré, bascule
en zone élevée dans les scénarios adverses~: la probabilité de franchir le seuil de décision
($P(\text{PD}\geq 0{,}5)$) devient substantielle, et la répartition par classe de risque montre
une migration nette vers la classe \textit{high}. Les scénarios extrêmes restitués — avec les
valeurs choquées qui les ont produits — permettent d'identifier précisément la combinaison de
chocs la plus pénalisante. Cette lecture transforme une décision binaire en une appréciation
nuancée de la \emph{robustesse} du dossier, directement exploitable en comité de crédit.

\subsection{Intégration à la plateforme}

Le moteur s'inscrit dans la même architecture que le reste de la plateforme~: il réutilise
l'artefact déployé (préprocesseur + modèle) chargé et mis en cache au démarrage, expose ses
résultats via des schémas Pydantic typés (\texttt{MonteCarloResponse}), et peut être invoqué
depuis le tableau de bord. Ses sorties — synthèse d'indicateurs, histogramme et scénarios
extrêmes — sont, comme les autres sorties calculées, susceptibles d'être mises en récit par
la couche LLM sous la même règle de séparation stricte des rôles~: aucune valeur numérique
n'est produite par le modèle de langage.

\section*{Conclusion}

Cette release finalise la plateforme en dotant les analystes d'un moteur de tests de stress
capable de mesurer la robustesse d'un dossier PME face à des chocs macro-économiques et
sectoriels. En réutilisant le modèle déployé de façon vectorisée, il reconstitue la distribution
complète de la PD sous stress et en extrait des indicateurs de risque extrême (VaR, Expected
Shortfall, probabilité de dépassement du seuil), offrant une lecture du risque bien plus riche
qu'une prédiction ponctuelle. Au-delà de la valeur métier, ce moteur illustre une dernière fois
les principes d'ingénierie qui structurent toute la plateforme~: réutilisation stricte de
l'artefact déployé pour garantir la cohérence avec la prédiction nominale, vectorisation pour la
performance, validation défensive des entrées, et reproductibilité par graine fixe au service de
l'auditabilité.

% ════════════════════════════════════════════════════════════
%  CONCLUSION GÉNÉRALE
% ════════════════════════════════════════════════════════════
\chapterunnumbered{Conclusion générale}

Ce projet de fin d'études, réalisé au sein de MPSoft, a conduit la conception et le
développement complets d'une plateforme d'aide à la décision pour le risque de crédit des PME,
articulée autour d'un principe directeur~: allier la puissance prédictive de l'apprentissage
automatique à une explicabilité et une traçabilité de niveau réglementaire, sans jamais
dessaisir l'analyste de la décision finale.

La démarche s'est déroulée en quatre \emph{releases} cohérentes. La première a établi la
compréhension métier du crédit PME et constitué un socle de données synthétique robuste,
piloté par un facteur de risque latent et enrichi de plus de quatre-vingt-dix variables
métier-cohérentes. La deuxième a construit le c\oe ur prédictif~: une chaîne de prétraitement
reproductible, une ingénierie de variables discriminantes, la comparaison de deux familles de
modèles ensemblistes (Random Forest et XGBoost) sur trois périmètres de variables (options~A,
B et~C), et le déploiement d'un XGBoost hybride (option~C) doté d'une couche d'explicabilité
SHAP globale et locale. La troisième a industrialisé ce modèle au sein d'un agent IA
orchestrateur, sous séparation stricte des rôles avec un LLM local chargé de la seule mise en
récit des résultats. La quatrième, enfin, a ajouté un moteur de tests de stress Monte Carlo
reconstituant la distribution de la PD sous chocs et ses indicateurs de risque extrême.

La solution livrée répond directement aux limites identifiées des processus d'instruction
existants~: elle réduit un traitement expert de plusieurs jours à une analyse automatisée de
quelques minutes, tout en préservant — et c'est là sa valeur distinctive — l'explicabilité de
chaque décision. Chaque prédiction est décomposable en contributions de variables, chaque
résultat est traçable et rattaché à une version de modèle, et l'architecture garantit qu'aucune
valeur numérique n'est produite par le LLM. La plateforme autonomise ainsi les analystes sans
automatiser le jugement du comité de crédit.

\medskip
\noindent\textbf{Apports techniques.} Au-delà des fonctionnalités, le projet a mis en \oe uvre
un ensemble de principes d'ingénierie qui font sa robustesse. La séparation stricte entre la
couche HTTP (\textit{routers}) et la logique métier (\textit{services}) rend chaque brique
réutilisable~: un même service de prédiction alimente à la fois la prédiction directe, la
simulation de scénarios, l'orchestration et le moteur de stress. La sérialisation conjointe du
préprocesseur et du modèle dans un artefact unique garantit une inférence strictement alignée
sur l'entraînement. Enfin, la conception défensive — validation non bloquante des entrées,
capture des pannes du LLM sans perte des résultats calculés, reproductibilité par graine fixe —
assure que la valeur métier est toujours délivrée et toujours auditable.

\medskip
\noindent\textbf{Limites.} Le principal facteur à relativiser tient à la nature
\emph{synthétique} du jeu de données~: la cible étant dérivée d'un score latent réutilisant des
signaux d'entrée, les performances observées (AUC $\approx$ 0{,}89) valident la cohérence
technique du pipeline mais ne préjugent pas d'un niveau de performance réaliste sur données
réelles. L'injection de réalisme (bruit, atténuation des variables dominantes, découpage
temporel) ramène d'ailleurs l'AUC à des niveaux plus crédibles ($\approx$ 0{,}80–0{,}82), ce qui
confirme que le score élevé initial reflétait surtout la cohérence interne du générateur. Par
ailleurs, les hyperparamètres ont été fixés à des valeurs éprouvées plutôt que recherchés
exhaustivement, et l'évaluation repose sur une partition \textit{train}/\textit{test} unique
plutôt que sur une validation croisée complète.

\medskip
\noindent\textbf{Perspectives.} Plusieurs prolongements se dégagent naturellement. Le plus
structurant est la \emph{migration vers les données réelles} de la banque~: ré-entraînement et
validation externe du pipeline sur les dossiers réels, calibration des probabilités et suivi de
la dérive du modèle dans le temps. La récupération des variables réellement prédictives présentes
dans le système d'information — sous-facteurs de notation, ratios financiers détaillés, historique
de paiement — conditionne la valeur d'un futur modèle déployable. À cela s'ajoutent une recherche
systématique d'hyperparamètres et une validation croisée temporelle, l'intégration de la PD dans
le calcul réglementaire des provisions pour risque de crédit, l'audit d'équité et de robustesse
des décisions, et enfin l'industrialisation (déploiement conteneurisé, supervision, montée à
l'échelle du portefeuille) et l'extension à d'autres segments ou clients de MPSoft.

% ════════════════════════════════════════════════════════════
%  WEBOGRAPHIE & BIBLIOGRAPHIE
% ════════════════════════════════════════════════════════════
\begin{thebibliography}{99}
\addcontentsline{toc}{chapter}{Webographie \& Bibliographie}

% ── Risque de crédit et scoring ──
\bibitem{altman1968} E.~I. Altman, ``Financial Ratios, Discriminant Analysis and the
  Prediction of Corporate Bankruptcy'', \emph{The Journal of Finance}, vol.~23, n\up{o}~4,
  p.~589--609, 1968.

\bibitem{thomas2002} L.~C. Thomas, D.~B. Edelman et J.~N. Crook, \emph{Credit Scoring and
  Its Applications}, Society for Industrial and Applied Mathematics (SIAM), Philadelphie, 2002.

\bibitem{hosmer2013} D.~W. Hosmer, S.~Lemeshow et R.~X. Sturdivant, \emph{Applied Logistic
  Regression}, 3\up{e} éd., Wiley, 2013.

\bibitem{lessmann2015} S.~Lessmann, B.~Baesens, H.-V. Seow et L.~C. Thomas, ``Benchmarking
  State-of-the-Art Classification Algorithms for Credit Scoring: An Update of Research'',
  \emph{European Journal of Operational Research}, vol.~247, n\up{o}~1, p.~124--136, 2015.

% ── Apprentissage automatique ──
\bibitem{breiman2001} L.~Breiman, ``Random Forests'', \emph{Machine Learning}, vol.~45,
  n\up{o}~1, p.~5--32, 2001.

\bibitem{friedman2001} J.~H. Friedman, ``Greedy Function Approximation: A Gradient Boosting
  Machine'', \emph{The Annals of Statistics}, vol.~29, n\up{o}~5, p.~1189--1232, 2001.

\bibitem{chen2016} T.~Chen et C.~Guestrin, ``XGBoost: A Scalable Tree Boosting System'',
  \emph{Proc. 22nd ACM SIGKDD Int. Conf. on Knowledge Discovery and Data Mining},
  p.~785--794, 2016.

\bibitem{pedregosa2011} F.~Pedregosa \textit{et al.}, ``Scikit-learn: Machine Learning in
  Python'', \emph{Journal of Machine Learning Research}, vol.~12, p.~2825--2830, 2011.

% ── Explicabilité (XAI) ──
\bibitem{lundberg2017} S.~M. Lundberg et S.-I. Lee, ``A Unified Approach to Interpreting Model
  Predictions'', \emph{Advances in Neural Information Processing Systems (NeurIPS)}, vol.~30,
  p.~4765--4774, 2017.

\bibitem{ribeiro2016} M.~T. Ribeiro, S.~Singh et C.~Guestrin, ``\,`Why Should I Trust You?':
  Explaining the Predictions of Any Classifier'', \emph{Proc. 22nd ACM SIGKDD Int. Conf. on
  Knowledge Discovery and Data Mining}, p.~1135--1144, 2016.

\bibitem{molnar2022} C.~Molnar, \emph{Interpretable Machine Learning: A Guide for Making Black
  Box Models Explainable}, 2\up{e} éd., 2022. \url{https://christophm.github.io/interpretable-ml-book/}

% ── Agents IA et grands modèles de langage ──
\bibitem{vaswani2017} A.~Vaswani \textit{et al.}, ``Attention Is All You Need'',
  \emph{Advances in Neural Information Processing Systems (NeurIPS)}, vol.~30, 2017.

\bibitem{brown2020} T.~B. Brown \textit{et al.}, ``Language Models Are Few-Shot Learners'',
  \emph{Advances in Neural Information Processing Systems (NeurIPS)}, vol.~33, p.~1877--1901, 2020.

\bibitem{wu2023bloomberggpt} S.~Wu \textit{et al.}, ``BloombergGPT: A Large Language Model for
  Finance'', \emph{arXiv preprint} arXiv:2303.17564, 2023.

% ── Tests de stress et Monte Carlo ──
\bibitem{glasserman2004} P.~Glasserman, \emph{Monte Carlo Methods in Financial Engineering},
  Springer, 2004.

\bibitem{bcbs2018} Basel Committee on Banking Supervision, \emph{Stress Testing Principles},
  Bank for International Settlements, 2018. \url{https://www.bis.org/bcbs/publ/d450.htm}

\bibitem{metropolis1949} N.~Metropolis et S.~Ulam, ``The Monte Carlo Method'', \emph{Journal
  of the American Statistical Association}, vol.~44, n\up{o}~247, p.~335--341, 1949.

% ── Outils et frameworks (webographie) ──
\bibitem{fastapi} S.~Ramírez, \emph{FastAPI --- Documentation officielle}, 2024.
  \url{https://fastapi.tiangolo.com/} (consulté le 01/05/2025).

\bibitem{pydantic} Pydantic, \emph{Pydantic --- Data Validation Using Python Type Hints}, 2024.
  \url{https://docs.pydantic.dev/} (consulté le 01/05/2025).

\bibitem{shapdoc} S.~M. Lundberg, \emph{SHAP --- SHapley Additive exPlanations, Documentation},
  2024. \url{https://shap.readthedocs.io/} (consulté le 01/05/2025).

\bibitem{xgboostdoc} XGBoost Developers, \emph{XGBoost Documentation}, 2024.
  \url{https://xgboost.readthedocs.io/} (consulté le 01/05/2025).

\bibitem{ollama} Ollama, \emph{Ollama --- Run Large Language Models Locally}, 2024.
  \url{https://ollama.com/} (consulté le 01/05/2025).

\end{thebibliography}

% ════════════════════════════════════════════════════════════
%  RÉSUMÉ ET ABSTRACT (à la fin du rapport)
% ════════════════════════════════════════════════════════════
\chapterunnumbered{Résumé}

Ce projet de fin d'études, réalisé au sein de MPSoft, vise à concevoir une
plateforme d'aide à la décision pour le risque de crédit des PME, basée sur un agent
d'intelligence artificielle explicable. La solution compare des modèles ensemblistes
(Random Forest et XGBoost) sur trois périmètres de variables et déploie un XGBoost hybride
pour estimer la probabilité de défaut, une couche d'explicabilité fondée sur SHAP, un agent
orchestrateur doté d'un module de raisonnement en langage naturel par LLM local, et des tests
de stress par simulation Monte Carlo. L'ensemble garantit l'auditabilité, la traçabilité et la
supervision humaine des décisions de crédit.

\bigskip
\noindent\textbf{Mots-clés :} Risque de crédit, Intelligence Artificielle, Explainable
AI, SHAP, Probabilité de Défaut, Agent IA, LLM, Monte Carlo, Scoring bancaire.

\chapterunnumbered{Abstract}

% [Abstract in English to write]
This graduation project, carried out at MPSoft, aims to design a decision-support
platform for SME credit risk based on an explainable AI agent. The solution compares
ensemble models (Random Forest and XGBoost) over three feature sets and deploys a hybrid
XGBoost model to estimate the probability of default, an explainability layer based on SHAP,
an orchestrating agent with a natural-language reasoning module powered by a local LLM, and
stress testing via Monte Carlo simulation. The platform ensures auditability, traceability, and
human-in-the-loop supervision of credit decisions.

\bigskip
\noindent\textbf{Keywords:} Credit Risk, Artificial Intelligence, Explainable AI,
SHAP, Probability of Default, AI Agent, LLM, Monte Carlo, Banking Scorecard.

% ════════════════════════════════════════════════════════════
%  TABLE DES FIGURES ET LISTE DES TABLEAUX (à la fin)
% ════════════════════════════════════════════════════════════
\clearpage
\renewcommand{\listfigurename}{TABLE DES FIGURES}
\listoffigures

\clearpage
\renewcommand{\listtablename}{LISTE DES TABLEAUX}
\listoftables

\end{document}
